%% ============================================================================
%% Causal4j -- ICDE 2027 Research Track (Regular paper)
%% IEEE conference format. Compile: pdflatex main -> bibtex main -> pdflatex x2
%% Template: IEEEtran (https://www.ieee.org/conferences_events/conferences/publishing/templates.html)
%% Length: <= 12 pages EXCLUDING references and the AI-content acknowledgement.
%% Single-blind: author names ARE shown at submission. ORCIDs required.
%% ============================================================================
\documentclass[conference]{IEEEtran}
\IEEEoverridecommandlockouts

\usepackage{cite}
\usepackage{amsmath,amssymb,amsfonts}
\usepackage{graphicx}
\usepackage{textcomp}
\usepackage{xcolor}
\usepackage{booktabs}
\usepackage{listings}
\usepackage{url}
\usepackage{tikz}
\usetikzlibrary{arrows.meta,positioning,fit,backgrounds,shapes.geometric,topaths,calc}

% --- lightweight DSL listing style (rename/extend as needed) ---
\lstdefinestyle{causaldsl}{
  basicstyle=\ttfamily\small,
  keywordstyle=\bfseries,
  columns=fullflexible,
  breaklines=true,
  keywords={DEFINE,MODEL,CAUSAL,DAG,VARIABLE,IDENTIFY,EFFECT,OF,ON,RETURN,
            ADJUSTMENT,SETS,ANCESTORS,DESCENDANTS,MEDIATORS,COLLIDERS,
            CONFOUNDERS,BACKDOOR,PATHS,BETWEEN,AND,ESTIMATE,INTERVENE,
            COMPARE,TO,COUNTERFACTUAL,GIVEN,WHAT,WOULD,HAVE,BEEN,IF,USE,RESET}
}

% --- editorial notes: remove before submission ---
\newcommand{\TODO}[1]{\textcolor{red}{[\textbf{TODO:} #1]}}
\newcommand{\NOTE}[1]{\textcolor{blue}{[#1]}}

\begin{document}

\title{Causal4j: A Declarative Language and In-DBMS Reachability
Index for Causal Structural Queries over Graph Databases}

% ORCIDs are REQUIRED for every author at submission time.
% NOTE: '&' is a LaTeX special character -- it MUST be escaped as '\&'
% (see "CNRS LIRIS \& IUF" below), otherwise compilation fails.
% ORCIDs are REQUIRED for every author at submission; add them before submitting.
\author{
\IEEEauthorblockN{Amedeo Pachera}
\IEEEauthorblockA{\textit{CNRS LIRIS} \\
\textit{Lyon 1 University}\\
Lyon, France \\
amedeo.pachera@univ-lyon1.fr}
\and
\IEEEauthorblockN{Kirollos Nayrouz}
\IEEEauthorblockA{\textit{CNRS LIRIS} \\
\textit{Lyon 1 University}\\
Lyon, France \\
kirollos.nayrouz@etu.univ-lyon1.fr}
\and
\IEEEauthorblockN{Angela Bonifati}
\IEEEauthorblockA{\textit{CNRS LIRIS \& IUF} \\
\textit{Lyon 1 University}\\
Lyon, France \\
angela.bonifati@univ-lyon1.fr}
\and
\IEEEauthorblockN{Andrea Mauri}
\IEEEauthorblockA{\textit{CNRS LIRIS} \\
\textit{Lyon 1 University}\\
Lyon, France \\
andrea.mauri@univ-lyon1.fr}
}

\maketitle

\begin{abstract}
\TODO{150--200 words. One paragraph. Problem (causal structural queries are
asked over graph-resident data but no DBMS treats them as a first-class,
indexed workload); approach (a declarative causal query language embedded in a
graph DBMS + a compiled reachability index: dense CSR adjacency, topological-layer
pruning, and precomputed ancestor/descendant bitsets); result (structural
queries drop from per-query $O(V{+}E)$ traversal to $O(V/64)$ bitset operations;
$N\times$ / $M\times$ speedups over in-DBMS traversal and external causal
libraries; sub-second index builds and incremental maintenance under updates).
Name datasets and the headline number.}
\end{abstract}

\begin{IEEEkeywords}
graph databases, query languages, reachability indexing, query processing,
causal DAGs, adjustment sets, incremental maintenance
\end{IEEEkeywords}

%% ----------------------------------------------------------------------------
% Architecture figure (referenced from the Introduction as Fig.~\ref{fig:arch}).
\begin{figure}[t]\centering
  \resizebox{\linewidth}{!}{\input{figures/architecture}}
  \caption{Causal4j architecture. Language-layer statements are parsed and
  compiled to Cypher writes and \texttt{causal.*} procedure calls; the plugin
  maintains a compiled in-memory index (CSR adjacency + ancestor/descendant
  bitsets) over the stored property graph and answers structural queries against
  it.}
  \label{fig:arch}
\end{figure}

%% ============================================================================
%% Section: Introduction
%% Included from main.tex via %% ============================================================================
%% Section: Introduction
%% Included from main.tex via %% ============================================================================
%% Section: Introduction
%% Included from main.tex via \input{sections/introduction}
%% Editorial notes use \TODO{} / \NOTE{} (defined in main.tex); remove before CR.
%% Numbers marked \TODO are NOT yet measured -- see ICDE2027-outline.md, sec. 13.
%% ============================================================================
\section{Introduction}\label{sec:intro}

Causal reasoning has become central to data-driven decision making across
epidemiology, economics, and the empirical sciences. Analysts increasingly ask
\emph{structural} questions of a causal model: \emph{which variables confound
the effect of a treatment $T$ on an outcome $Y$?}, \emph{does a valid adjustment
set exist, and what is it?}, or \emph{is $X$ independent of $Y$ once we condition
on $Z$?} Each of these is answered not from data alone but from the topology of a
\emph{causal directed acyclic graph} (DAG), whose nodes are variables and whose
edges encode direct causal influence~\cite{pearl2009causality}. Such DAGs are no
longer small, hand-drawn diagrams: they are large, continually revised, and
increasingly stored and versioned inside graph databases alongside the domain
data they describe.

Despite this, causal analysis today is largely divorced from the systems that
store the graph. The prevailing practice is to \emph{export} the DAG into a
standalone library---\textsc{DoWhy}, \textsc{dagitty}, \textsc{pgmpy}, or
\textsc{causal-learn}---compute the answer in a separate process, and discard the
in-memory representation afterwards. This round-trip forfeits precisely the
properties a database is built to provide: durable and concurrent access to a
single authoritative model, declarative querying, integration with co-located
data, and---critically---\emph{incremental} response to a changing graph. When an
edge is added or a variable is reclassified as latent, the entire model must be
re-extracted and every downstream index rebuilt from scratch. No graph DBMS
today treats causal structural queries as a first-class, \emph{indexed} workload.

Closing this gap is not merely a matter of wrapping a library behind a stored
procedure. First, the operators analysts need are not all expressible as plain
graph reachability: \emph{d-separation} and \emph{adjustment-set enumeration}
depend on the conditioning set and on edge orientation, and require
conditioning-aware traversal (the Bayes-Ball procedure) rather than a transitive
closure. Second, answered naively---by traversing the stored graph on every
call---each query costs $O(V{+}E)$, and identification tasks issue many such
queries as they enumerate and test candidate adjustment sets, so per-query
traversal does not scale to large DAGs or high query volume. Third, the index
that accelerates these queries must survive updates: a causal graph under active
curation cannot afford a full recompilation after every edit.

\NOTE{Optional 4th challenge sentence: expressing the operators in a language
that is close to how domain experts phrase causal questions, yet compiles to
physical operators a query engine can execute.}

We present \textbf{Causal4j}, a system that makes causal structural queries a
native, indexed workload inside a property-graph DBMS (Neo4j). Causal4j has
three parts, shown in Fig.~\ref{fig:arch}. (i) A \emph{declarative causal query
language} whose operators---ancestors and descendants, mediators, colliders and
confounders, backdoor paths, d-separation, and adjustment-set
identification---read like the questions analysts already ask, and whose
statements compile to physical operators over the graph store. (ii) A
\emph{compiled causal index}: on compilation the DAG stored as
$(:$\textsf{Variable}$)\text{-}[:$\textsf{CAUSES}$]\!\rightarrow\!(:$\textsf{Variable}$)$
is snapshotted into a dense, integer-indexed layout (CSR adjacency, a
topological order and per-node layer) augmented with precomputed
\emph{ancestor} and \emph{descendant} bitsets; structural queries then reduce
from $O(V{+}E)$ traversals to $O(V/64)$ word operations, and d-separation and
adjustment enumeration are cast as bitset intersections. (iii) \emph{Incremental
maintenance} of that index under edge and variable updates, so a mutating causal
graph is kept queryable without full recompilation.

\NOTE{Insert Fig.~1 (architecture) here -- provided by figures/architecture.tex.}

This paper makes the following contributions:
\begin{itemize}
  \item A \textbf{declarative causal query language}, embedded in a graph DBMS,
  covering structural queries, identification (adjustment sets, backdoor
  criterion), and a quantitative layer (interventions, counterfactuals), with a
  defined compilation from language operators to physical graph operators
  (\S\ref{sec:lang}).
  \item A \textbf{compiled reachability index for causal DAGs}---dense CSR
  adjacency with topological-layer pruning and precomputed ancestor/descendant
  bitsets---that reduces structural queries from $O(V{+}E)$ to $O(V/64)$ and
  supports bitset Bayes-Ball d-separation and adjustment-set enumeration
  (\S\ref{sec:index}).
  \item An \textbf{incremental maintenance} scheme that updates the index under
  edge and variable mutations, bounding work to the affected region rather than
  recompiling the graph (\S\ref{sec:incr}).
  \item An \textbf{empirical study} on standard Bayesian-network DAGs and
  synthetic scale-free DAGs, comparing Causal4j against in-DBMS Cypher traversal,
  Neo4j GDS, and external causal libraries on latency, throughput, index build
  time, memory, and update cost, released as a reproducible artifact
  (\S\ref{sec:eval}).
\end{itemize}

\NOTE{Results preview -- fill once the eval in \S\ref{sec:eval} exists. Keep it
one or two sentences with the headline numbers.}
On DAGs ranging from \TODO{$V_{\min}$} to \TODO{$V_{\max}$} variables, Causal4j
answers d-separation and adjustment-set queries \TODO{$N\times$} faster than
in-DBMS Cypher traversal and \TODO{$M\times$} faster than round-tripping to an
external library, builds its index in \TODO{under X\,ms}, and applies an edge
update in \TODO{Y}, versus \TODO{Z} for a full recompilation.

\emph{Scope note for reviewers.} This is a data-engineering contribution: a
declarative query language embedded in a graph DBMS, a physical reachability index
for causal DAGs, query processing over that index, and its incremental
maintenance under updates. The structural-causal-model semantics (d-separation,
the backdoor criterion, adjustment sets) are the \emph{workload} the engine
supports, drawn from established causal theory~\cite{pearl2009causality}; they are
not claimed as a contribution. What is new is treating that workload as a
first-class, indexed, incrementally maintained query workload \emph{inside} the
store, rather than in an external library.

The remainder of the paper is organized as follows. \S\ref{sec:bg} defines
causal DAGs, the query catalog, and the in-DBMS problem. \S\ref{sec:lang}
presents the language and its compilation; \S\ref{sec:compile} the compilation
of the causal snapshot; \S\ref{sec:index} the index and query processing;
\S\ref{sec:incr} incremental maintenance; and \S\ref{sec:quant} the quantitative
layer. \S\ref{sec:eval} reports the evaluation, \S\ref{sec:related} surveys
related work, and \S\ref{sec:concl} concludes.

%% Editorial notes use \TODO{} / \NOTE{} (defined in main.tex); remove before CR.
%% Numbers marked \TODO are NOT yet measured -- see ICDE2027-outline.md, sec. 13.
%% ============================================================================
\section{Introduction}\label{sec:intro}

Causal reasoning has become central to data-driven decision making across
epidemiology, economics, and the empirical sciences. Analysts increasingly ask
\emph{structural} questions of a causal model: \emph{which variables confound
the effect of a treatment $T$ on an outcome $Y$?}, \emph{does a valid adjustment
set exist, and what is it?}, or \emph{is $X$ independent of $Y$ once we condition
on $Z$?} Each of these is answered not from data alone but from the topology of a
\emph{causal directed acyclic graph} (DAG), whose nodes are variables and whose
edges encode direct causal influence~\cite{pearl2009causality}. Such DAGs are no
longer small, hand-drawn diagrams: they are large, continually revised, and
increasingly stored and versioned inside graph databases alongside the domain
data they describe.

Despite this, causal analysis today is largely divorced from the systems that
store the graph. The prevailing practice is to \emph{export} the DAG into a
standalone library---\textsc{DoWhy}, \textsc{dagitty}, \textsc{pgmpy}, or
\textsc{causal-learn}---compute the answer in a separate process, and discard the
in-memory representation afterwards. This round-trip forfeits precisely the
properties a database is built to provide: durable and concurrent access to a
single authoritative model, declarative querying, integration with co-located
data, and---critically---\emph{incremental} response to a changing graph. When an
edge is added or a variable is reclassified as latent, the entire model must be
re-extracted and every downstream index rebuilt from scratch. No graph DBMS
today treats causal structural queries as a first-class, \emph{indexed} workload.

Closing this gap is not merely a matter of wrapping a library behind a stored
procedure. First, the operators analysts need are not all expressible as plain
graph reachability: \emph{d-separation} and \emph{adjustment-set enumeration}
depend on the conditioning set and on edge orientation, and require
conditioning-aware traversal (the Bayes-Ball procedure) rather than a transitive
closure. Second, answered naively---by traversing the stored graph on every
call---each query costs $O(V{+}E)$, and identification tasks issue many such
queries as they enumerate and test candidate adjustment sets, so per-query
traversal does not scale to large DAGs or high query volume. Third, the index
that accelerates these queries must survive updates: a causal graph under active
curation cannot afford a full recompilation after every edit.

\NOTE{Optional 4th challenge sentence: expressing the operators in a language
that is close to how domain experts phrase causal questions, yet compiles to
physical operators a query engine can execute.}

We present \textbf{Causal4j}, a system that makes causal structural queries a
native, indexed workload inside a property-graph DBMS (Neo4j). Causal4j has
three parts, shown in Fig.~\ref{fig:arch}. (i) A \emph{declarative causal query
language} whose operators---ancestors and descendants, mediators, colliders and
confounders, backdoor paths, d-separation, and adjustment-set
identification---read like the questions analysts already ask, and whose
statements compile to physical operators over the graph store. (ii) A
\emph{compiled causal index}: on compilation the DAG stored as
$(:$\textsf{Variable}$)\text{-}[:$\textsf{CAUSES}$]\!\rightarrow\!(:$\textsf{Variable}$)$
is snapshotted into a dense, integer-indexed layout (CSR adjacency, a
topological order and per-node layer) augmented with precomputed
\emph{ancestor} and \emph{descendant} bitsets; structural queries then reduce
from $O(V{+}E)$ traversals to $O(V/64)$ word operations, and d-separation and
adjustment enumeration are cast as bitset intersections. (iii) \emph{Incremental
maintenance} of that index under edge and variable updates, so a mutating causal
graph is kept queryable without full recompilation.

\NOTE{Insert Fig.~1 (architecture) here -- provided by figures/architecture.tex.}

This paper makes the following contributions:
\begin{itemize}
  \item A \textbf{declarative causal query language}, embedded in a graph DBMS,
  covering structural queries, identification (adjustment sets, backdoor
  criterion), and a quantitative layer (interventions, counterfactuals), with a
  defined compilation from language operators to physical graph operators
  (\S\ref{sec:lang}).
  \item A \textbf{compiled reachability index for causal DAGs}---dense CSR
  adjacency with topological-layer pruning and precomputed ancestor/descendant
  bitsets---that reduces structural queries from $O(V{+}E)$ to $O(V/64)$ and
  supports bitset Bayes-Ball d-separation and adjustment-set enumeration
  (\S\ref{sec:index}).
  \item An \textbf{incremental maintenance} scheme that updates the index under
  edge and variable mutations, bounding work to the affected region rather than
  recompiling the graph (\S\ref{sec:incr}).
  \item An \textbf{empirical study} on standard Bayesian-network DAGs and
  synthetic scale-free DAGs, comparing Causal4j against in-DBMS Cypher traversal,
  Neo4j GDS, and external causal libraries on latency, throughput, index build
  time, memory, and update cost, released as a reproducible artifact
  (\S\ref{sec:eval}).
\end{itemize}

\NOTE{Results preview -- fill once the eval in \S\ref{sec:eval} exists. Keep it
one or two sentences with the headline numbers.}
On DAGs ranging from \TODO{$V_{\min}$} to \TODO{$V_{\max}$} variables, Causal4j
answers d-separation and adjustment-set queries \TODO{$N\times$} faster than
in-DBMS Cypher traversal and \TODO{$M\times$} faster than round-tripping to an
external library, builds its index in \TODO{under X\,ms}, and applies an edge
update in \TODO{Y}, versus \TODO{Z} for a full recompilation.

\emph{Scope note for reviewers.} This is a data-engineering contribution: a
declarative query language embedded in a graph DBMS, a physical reachability index
for causal DAGs, query processing over that index, and its incremental
maintenance under updates. The structural-causal-model semantics (d-separation,
the backdoor criterion, adjustment sets) are the \emph{workload} the engine
supports, drawn from established causal theory~\cite{pearl2009causality}; they are
not claimed as a contribution. What is new is treating that workload as a
first-class, indexed, incrementally maintained query workload \emph{inside} the
store, rather than in an external library.

The remainder of the paper is organized as follows. \S\ref{sec:bg} defines
causal DAGs, the query catalog, and the in-DBMS problem. \S\ref{sec:lang}
presents the language and its compilation; \S\ref{sec:compile} the compilation
of the causal snapshot; \S\ref{sec:index} the index and query processing;
\S\ref{sec:incr} incremental maintenance; and \S\ref{sec:quant} the quantitative
layer. \S\ref{sec:eval} reports the evaluation, \S\ref{sec:related} surveys
related work, and \S\ref{sec:concl} concludes.

%% Editorial notes use \TODO{} / \NOTE{} (defined in main.tex); remove before CR.
%% Numbers marked \TODO are NOT yet measured -- see ICDE2027-outline.md, sec. 13.
%% ============================================================================
\section{Introduction}\label{sec:intro}

Causal reasoning has become central to data-driven decision making across
epidemiology, economics, and the empirical sciences. Analysts increasingly ask
\emph{structural} questions of a causal model: \emph{which variables confound
the effect of a treatment $T$ on an outcome $Y$?}, \emph{does a valid adjustment
set exist, and what is it?}, or \emph{is $X$ independent of $Y$ once we condition
on $Z$?} Each of these is answered not from data alone but from the topology of a
\emph{causal directed acyclic graph} (DAG), whose nodes are variables and whose
edges encode direct causal influence~\cite{pearl2009causality}. Such DAGs are no
longer small, hand-drawn diagrams: they are large, continually revised, and
increasingly stored and versioned inside graph databases alongside the domain
data they describe.

Despite this, causal analysis today is largely divorced from the systems that
store the graph. The prevailing practice is to \emph{export} the DAG into a
standalone library---\textsc{DoWhy}, \textsc{dagitty}, \textsc{pgmpy}, or
\textsc{causal-learn}---compute the answer in a separate process, and discard the
in-memory representation afterwards. This round-trip forfeits precisely the
properties a database is built to provide: durable and concurrent access to a
single authoritative model, declarative querying, integration with co-located
data, and---critically---\emph{incremental} response to a changing graph. When an
edge is added or a variable is reclassified as latent, the entire model must be
re-extracted and every downstream index rebuilt from scratch. No graph DBMS
today treats causal structural queries as a first-class, \emph{indexed} workload.

Closing this gap is not merely a matter of wrapping a library behind a stored
procedure. First, the operators analysts need are not all expressible as plain
graph reachability: \emph{d-separation} and \emph{adjustment-set enumeration}
depend on the conditioning set and on edge orientation, and require
conditioning-aware traversal (the Bayes-Ball procedure) rather than a transitive
closure. Second, answered naively---by traversing the stored graph on every
call---each query costs $O(V{+}E)$, and identification tasks issue many such
queries as they enumerate and test candidate adjustment sets, so per-query
traversal does not scale to large DAGs or high query volume. Third, the index
that accelerates these queries must survive updates: a causal graph under active
curation cannot afford a full recompilation after every edit.

\NOTE{Optional 4th challenge sentence: expressing the operators in a language
that is close to how domain experts phrase causal questions, yet compiles to
physical operators a query engine can execute.}

We present \textbf{Causal4j}, a system that makes causal structural queries a
native, indexed workload inside a property-graph DBMS (Neo4j). Causal4j has
three parts, shown in Fig.~\ref{fig:arch}. (i) A \emph{declarative causal query
language} whose operators---ancestors and descendants, mediators, colliders and
confounders, backdoor paths, d-separation, and adjustment-set
identification---read like the questions analysts already ask, and whose
statements compile to physical operators over the graph store. (ii) A
\emph{compiled causal index}: on compilation the DAG stored as
$(:$\textsf{Variable}$)\text{-}[:$\textsf{CAUSES}$]\!\rightarrow\!(:$\textsf{Variable}$)$
is snapshotted into a dense, integer-indexed layout (CSR adjacency, a
topological order and per-node layer) augmented with precomputed
\emph{ancestor} and \emph{descendant} bitsets; structural queries then reduce
from $O(V{+}E)$ traversals to $O(V/64)$ word operations, and d-separation and
adjustment enumeration are cast as bitset intersections. (iii) \emph{Incremental
maintenance} of that index under edge and variable updates, so a mutating causal
graph is kept queryable without full recompilation.

\NOTE{Insert Fig.~1 (architecture) here -- provided by figures/architecture.tex.}

This paper makes the following contributions:
\begin{itemize}
  \item A \textbf{declarative causal query language}, embedded in a graph DBMS,
  covering structural queries, identification (adjustment sets, backdoor
  criterion), and a quantitative layer (interventions, counterfactuals), with a
  defined compilation from language operators to physical graph operators
  (\S\ref{sec:lang}).
  \item A \textbf{compiled reachability index for causal DAGs}---dense CSR
  adjacency with topological-layer pruning and precomputed ancestor/descendant
  bitsets---that reduces structural queries from $O(V{+}E)$ to $O(V/64)$ and
  supports bitset Bayes-Ball d-separation and adjustment-set enumeration
  (\S\ref{sec:index}).
  \item An \textbf{incremental maintenance} scheme that updates the index under
  edge and variable mutations, bounding work to the affected region rather than
  recompiling the graph (\S\ref{sec:incr}).
  \item An \textbf{empirical study} on standard Bayesian-network DAGs and
  synthetic scale-free DAGs, comparing Causal4j against in-DBMS Cypher traversal,
  Neo4j GDS, and external causal libraries on latency, throughput, index build
  time, memory, and update cost, released as a reproducible artifact
  (\S\ref{sec:eval}).
\end{itemize}

\NOTE{Results preview -- fill once the eval in \S\ref{sec:eval} exists. Keep it
one or two sentences with the headline numbers.}
On DAGs ranging from \TODO{$V_{\min}$} to \TODO{$V_{\max}$} variables, Causal4j
answers d-separation and adjustment-set queries \TODO{$N\times$} faster than
in-DBMS Cypher traversal and \TODO{$M\times$} faster than round-tripping to an
external library, builds its index in \TODO{under X\,ms}, and applies an edge
update in \TODO{Y}, versus \TODO{Z} for a full recompilation.

\emph{Scope note for reviewers.} This is a data-engineering contribution: a
declarative query language embedded in a graph DBMS, a physical reachability index
for causal DAGs, query processing over that index, and its incremental
maintenance under updates. The structural-causal-model semantics (d-separation,
the backdoor criterion, adjustment sets) are the \emph{workload} the engine
supports, drawn from established causal theory~\cite{pearl2009causality}; they are
not claimed as a contribution. What is new is treating that workload as a
first-class, indexed, incrementally maintained query workload \emph{inside} the
store, rather than in an external library.

The remainder of the paper is organized as follows. \S\ref{sec:bg} defines
causal DAGs, the query catalog, and the in-DBMS problem. \S\ref{sec:lang}
presents the language and its compilation; \S\ref{sec:compile} the compilation
of the causal snapshot; \S\ref{sec:index} the index and query processing;
\S\ref{sec:incr} incremental maintenance; and \S\ref{sec:quant} the quantitative
layer. \S\ref{sec:eval} reports the evaluation, \S\ref{sec:related} surveys
related work, and \S\ref{sec:concl} concludes.


%% ----------------------------------------------------------------------------
% Running-example figure (referenced from Background as Fig.~\ref{fig:covid}).
\begin{figure}[t]\centering
  \resizebox{\linewidth}{!}{\input{figures/running-example}}
  \caption{A fragment of the COVID-19 causal DAG~\cite{piccininni2020covid}
  (full model: $25$ variables, $76$ edges). Treatment $T=\textsf{interventions}$
  and outcome $Y=\textsf{reported\_cases}$ are shaded; \textsf{covid19\_burden}
  is a confounder (dotted); dashed nodes are latent. The confounder opens a
  backdoor path $T\!\leftarrow\!\textsf{covid19\_burden}\!\rightarrow\!
  \textsf{exposure\_sarscov2}\!\rightarrow\!\textsf{onset\_covid19}\!\rightarrow\!Y$
  (bold), blocked by adjusting for $\{\textsf{covid19\_burden}\}$.}
  \label{fig:covid}
\end{figure}

%% ============================================================================
%% Section: Background and Problem
%% Included from main.tex via %% ============================================================================
%% Section: Background and Problem
%% Included from main.tex via %% ============================================================================
%% Section: Background and Problem
%% Included from main.tex via \input{sections/background}
%% ============================================================================
\section{Background and Problem}\label{sec:bg}

\subsection{Causal DAGs and their property-graph encoding}\label{sec:bg:dag}

A \emph{causal DAG} is a directed acyclic graph $G=(V,E)$ whose vertices $V$ are
random variables and whose edges $E$ encode direct causal influence: an edge
$X\!\rightarrow\!Y$ asserts that $X$ is a direct cause of $Y$ relative to the
other modeled variables~\cite{pearl2009causality}. Variables are either
\emph{observed} or \emph{latent} (unmeasured); latent variables are retained
because they change which conditional independencies hold and which causal
effects are identifiable, even though they are never assigned data. We write
$\mathrm{pa}(X)$, $\mathrm{ch}(X)$, $\mathrm{an}(X)$, and $\mathrm{de}(X)$ for
the parents, children, ancestors, and descendants of $X$; ancestors and
descendants are the transitive closures of parents and children and exclude $X$
itself.

Causal4j stores the DAG as a native property graph. Each variable is a node
\[
  (:\textsf{Variable}\;\{\textsf{id},\ \textsf{name},\ \textsf{observed},\ \textsf{latent}\})
\]
and each direct cause is a relationship
\[
  (:\textsf{Variable})\text{-}[:\textsf{CAUSES}\;\{\textsf{active},\ \textsf{weight}\}]\!\rightarrow\!(:\textsf{Variable}).
\]
The \textsf{id} property is the canonical key; \textsf{active} lets a curator
toggle an edge without deleting it, and only edges with $\textsf{active}=\text{true}$
participate in a model; \textsf{weight} carries an optional structural
coefficient used by the quantitative layer (\S\ref{sec:quant}) and ignored by the
structural queries. This encoding keeps the causal model \emph{in} the database,
co-located and jointly queryable with any other graph data, rather than in a
transient in-memory object owned by an external tool.

\subsection{The causal query catalog}\label{sec:bg:catalog}

Analysts interrogate the DAG through a small, well-established
catalog of structural queries. We group them by the reasoning they support.

\paragraph{Reachability queries} return $\mathrm{pa}$, $\mathrm{ch}$,
$\mathrm{an}$, and $\mathrm{de}$ of a variable, and test whether a directed path
$X\rightarrow\!\cdots\!\rightarrow Y$ exists. They underpin every other query.

\paragraph{Path-role queries} classify intermediate variables on the paths
between a treatment $T$ and an outcome $Y$. A \emph{mediator} lies on a directed
path $T\rightarrow\!\cdots\!\rightarrow Y$; a \emph{confounder} is a common cause,
i.e.\ a variable in $\mathrm{an}(T)\cap\mathrm{an}(Y)$; a \emph{collider} on a
path is a variable with two incoming path edges ($\rightarrow Z \leftarrow$).

\paragraph{Independence queries} test \emph{d-separation}. A path between $X$ and
$Y$ is \emph{blocked} by a conditioning set $Z$ if it contains a non-collider in
$Z$ or a collider none of whose descendants is in $Z$; $X$ and $Y$ are
d-separated given $Z$, written $X \!\perp\! Y \mid Z$, when \emph{every} path
between them is blocked~\cite{pearl2009causality}. D-separation is the graphical
criterion for conditional independence and is \emph{not} plain reachability: a
variable added to $Z$ can block one path while unblocking another through a
collider, so the test must be conditioning-aware.

\paragraph{Identification queries} decide whether the causal effect of $T$ on $Y$
can be estimated from observational data and, if so, return an \emph{adjustment
set}. A \emph{backdoor path} from $T$ to $Y$ is a path with an edge into $T$; a
set $Z$ satisfies the \emph{backdoor criterion} if it blocks every backdoor path
and contains no descendant of $T$. Enumerating minimal valid adjustment sets
issues many d-separation tests, so identification is the most traversal-intensive
query in the catalog.

Table~\ref{tab:operators} (\S\ref{sec:lang}) maps each of these to a Causal4j
language operator and the physical operator that answers it.

\subsection{Running example}\label{sec:bg:example}

We use throughout the COVID-19 causal DAG of Piccininni et
al.~\cite{piccininni2020covid}, comprising $25$ variables and $76$ directed
edges; Fig.~\ref{fig:covid} shows the fragment relevant to one query. Consider
estimating the effect of \textsf{interventions} ($T$) on
\textsf{reported\_cases} ($Y$). The causal (front-door) influence flows along
directed paths such as
$\textsf{interventions}\!\rightarrow\!\textsf{exposure\_sarscov2}\!\rightarrow\!
\textsf{onset\_covid19}\!\rightarrow\!\textsf{reported\_cases}$, whose interior
variables are mediators---and here happen to be \emph{latent}, which the engine
must nonetheless reason about. \textsf{covid19\_burden} is a common cause of both
$T$ and $Y$ (it causes \textsf{interventions} and also
\textsf{exposure\_sarscov2}), so it is a confounder that opens the backdoor path
$\textsf{interventions}\!\leftarrow\!\textsf{covid19\_burden}\!\rightarrow\!
\textsf{exposure\_sarscov2}\!\rightarrow\!\cdots\!\rightarrow\!\textsf{reported\_cases}$.
Conditioning on $\{\textsf{covid19\_burden}\}$ blocks this path and satisfies the
backdoor criterion, so it is a valid adjustment set for $T\!\rightarrow\!Y$. Every
claim in this paragraph is the output of a query in the catalog above.

\subsection{Problem statement}\label{sec:bg:problem}

We target the following data-engineering problem: \emph{support the causal query
catalog as a first-class workload inside a graph DBMS}. Concretely, a solution
must meet four requirements.
\begin{description}
  \item[R1 (Expressiveness).] A declarative interface in which each catalog query
  is a single statement, close to how analysts phrase causal questions, that
  compiles to physical operators over the graph store (\S\ref{sec:lang}).
  \item[R2 (Efficiency).] Structural queries answered without re-traversing the
  stored graph on every call, so that latency is sublinear in graph size and the
  system sustains the many independence tests an identification query issues
  (\S\ref{sec:index}).
  \item[R3 (Correctness).] Results faithful to the graphical semantics above,
  in particular a conditioning-aware d-separation test and a sound backdoor
  criterion (\S\ref{sec:index}, validated in \S\ref{sec:eval}).
  \item[R4 (Maintainability).] The accelerating index kept consistent under edge
  and variable updates without full recomputation, since curated causal graphs
  change continually (\S\ref{sec:incr}).
\end{description}
None of these is served by exporting the DAG to an external library: R2--R4 in
particular are database concerns---indexing, query processing, and incremental
maintenance---which is why we address them inside the DBMS.

%% ============================================================================
\section{Background and Problem}\label{sec:bg}

\subsection{Causal DAGs and their property-graph encoding}\label{sec:bg:dag}

A \emph{causal DAG} is a directed acyclic graph $G=(V,E)$ whose vertices $V$ are
random variables and whose edges $E$ encode direct causal influence: an edge
$X\!\rightarrow\!Y$ asserts that $X$ is a direct cause of $Y$ relative to the
other modeled variables~\cite{pearl2009causality}. Variables are either
\emph{observed} or \emph{latent} (unmeasured); latent variables are retained
because they change which conditional independencies hold and which causal
effects are identifiable, even though they are never assigned data. We write
$\mathrm{pa}(X)$, $\mathrm{ch}(X)$, $\mathrm{an}(X)$, and $\mathrm{de}(X)$ for
the parents, children, ancestors, and descendants of $X$; ancestors and
descendants are the transitive closures of parents and children and exclude $X$
itself.

Causal4j stores the DAG as a native property graph. Each variable is a node
\[
  (:\textsf{Variable}\;\{\textsf{id},\ \textsf{name},\ \textsf{observed},\ \textsf{latent}\})
\]
and each direct cause is a relationship
\[
  (:\textsf{Variable})\text{-}[:\textsf{CAUSES}\;\{\textsf{active},\ \textsf{weight}\}]\!\rightarrow\!(:\textsf{Variable}).
\]
The \textsf{id} property is the canonical key; \textsf{active} lets a curator
toggle an edge without deleting it, and only edges with $\textsf{active}=\text{true}$
participate in a model; \textsf{weight} carries an optional structural
coefficient used by the quantitative layer (\S\ref{sec:quant}) and ignored by the
structural queries. This encoding keeps the causal model \emph{in} the database,
co-located and jointly queryable with any other graph data, rather than in a
transient in-memory object owned by an external tool.

\subsection{The causal query catalog}\label{sec:bg:catalog}

Analysts interrogate the DAG through a small, well-established
catalog of structural queries. We group them by the reasoning they support.

\paragraph{Reachability queries} return $\mathrm{pa}$, $\mathrm{ch}$,
$\mathrm{an}$, and $\mathrm{de}$ of a variable, and test whether a directed path
$X\rightarrow\!\cdots\!\rightarrow Y$ exists. They underpin every other query.

\paragraph{Path-role queries} classify intermediate variables on the paths
between a treatment $T$ and an outcome $Y$. A \emph{mediator} lies on a directed
path $T\rightarrow\!\cdots\!\rightarrow Y$; a \emph{confounder} is a common cause,
i.e.\ a variable in $\mathrm{an}(T)\cap\mathrm{an}(Y)$; a \emph{collider} on a
path is a variable with two incoming path edges ($\rightarrow Z \leftarrow$).

\paragraph{Independence queries} test \emph{d-separation}. A path between $X$ and
$Y$ is \emph{blocked} by a conditioning set $Z$ if it contains a non-collider in
$Z$ or a collider none of whose descendants is in $Z$; $X$ and $Y$ are
d-separated given $Z$, written $X \!\perp\! Y \mid Z$, when \emph{every} path
between them is blocked~\cite{pearl2009causality}. D-separation is the graphical
criterion for conditional independence and is \emph{not} plain reachability: a
variable added to $Z$ can block one path while unblocking another through a
collider, so the test must be conditioning-aware.

\paragraph{Identification queries} decide whether the causal effect of $T$ on $Y$
can be estimated from observational data and, if so, return an \emph{adjustment
set}. A \emph{backdoor path} from $T$ to $Y$ is a path with an edge into $T$; a
set $Z$ satisfies the \emph{backdoor criterion} if it blocks every backdoor path
and contains no descendant of $T$. Enumerating minimal valid adjustment sets
issues many d-separation tests, so identification is the most traversal-intensive
query in the catalog.

Table~\ref{tab:operators} (\S\ref{sec:lang}) maps each of these to a Causal4j
language operator and the physical operator that answers it.

\subsection{Running example}\label{sec:bg:example}

We use throughout the COVID-19 causal DAG of Piccininni et
al.~\cite{piccininni2020covid}, comprising $25$ variables and $76$ directed
edges; Fig.~\ref{fig:covid} shows the fragment relevant to one query. Consider
estimating the effect of \textsf{interventions} ($T$) on
\textsf{reported\_cases} ($Y$). The causal (front-door) influence flows along
directed paths such as
$\textsf{interventions}\!\rightarrow\!\textsf{exposure\_sarscov2}\!\rightarrow\!
\textsf{onset\_covid19}\!\rightarrow\!\textsf{reported\_cases}$, whose interior
variables are mediators---and here happen to be \emph{latent}, which the engine
must nonetheless reason about. \textsf{covid19\_burden} is a common cause of both
$T$ and $Y$ (it causes \textsf{interventions} and also
\textsf{exposure\_sarscov2}), so it is a confounder that opens the backdoor path
$\textsf{interventions}\!\leftarrow\!\textsf{covid19\_burden}\!\rightarrow\!
\textsf{exposure\_sarscov2}\!\rightarrow\!\cdots\!\rightarrow\!\textsf{reported\_cases}$.
Conditioning on $\{\textsf{covid19\_burden}\}$ blocks this path and satisfies the
backdoor criterion, so it is a valid adjustment set for $T\!\rightarrow\!Y$. Every
claim in this paragraph is the output of a query in the catalog above.

\subsection{Problem statement}\label{sec:bg:problem}

We target the following data-engineering problem: \emph{support the causal query
catalog as a first-class workload inside a graph DBMS}. Concretely, a solution
must meet four requirements.
\begin{description}
  \item[R1 (Expressiveness).] A declarative interface in which each catalog query
  is a single statement, close to how analysts phrase causal questions, that
  compiles to physical operators over the graph store (\S\ref{sec:lang}).
  \item[R2 (Efficiency).] Structural queries answered without re-traversing the
  stored graph on every call, so that latency is sublinear in graph size and the
  system sustains the many independence tests an identification query issues
  (\S\ref{sec:index}).
  \item[R3 (Correctness).] Results faithful to the graphical semantics above,
  in particular a conditioning-aware d-separation test and a sound backdoor
  criterion (\S\ref{sec:index}, validated in \S\ref{sec:eval}).
  \item[R4 (Maintainability).] The accelerating index kept consistent under edge
  and variable updates without full recomputation, since curated causal graphs
  change continually (\S\ref{sec:incr}).
\end{description}
None of these is served by exporting the DAG to an external library: R2--R4 in
particular are database concerns---indexing, query processing, and incremental
maintenance---which is why we address them inside the DBMS.

%% ============================================================================
\section{Background and Problem}\label{sec:bg}

\subsection{Causal DAGs and their property-graph encoding}\label{sec:bg:dag}

A \emph{causal DAG} is a directed acyclic graph $G=(V,E)$ whose vertices $V$ are
random variables and whose edges $E$ encode direct causal influence: an edge
$X\!\rightarrow\!Y$ asserts that $X$ is a direct cause of $Y$ relative to the
other modeled variables~\cite{pearl2009causality}. Variables are either
\emph{observed} or \emph{latent} (unmeasured); latent variables are retained
because they change which conditional independencies hold and which causal
effects are identifiable, even though they are never assigned data. We write
$\mathrm{pa}(X)$, $\mathrm{ch}(X)$, $\mathrm{an}(X)$, and $\mathrm{de}(X)$ for
the parents, children, ancestors, and descendants of $X$; ancestors and
descendants are the transitive closures of parents and children and exclude $X$
itself.

Causal4j stores the DAG as a native property graph. Each variable is a node
\[
  (:\textsf{Variable}\;\{\textsf{id},\ \textsf{name},\ \textsf{observed},\ \textsf{latent}\})
\]
and each direct cause is a relationship
\[
  (:\textsf{Variable})\text{-}[:\textsf{CAUSES}\;\{\textsf{active},\ \textsf{weight}\}]\!\rightarrow\!(:\textsf{Variable}).
\]
The \textsf{id} property is the canonical key; \textsf{active} lets a curator
toggle an edge without deleting it, and only edges with $\textsf{active}=\text{true}$
participate in a model; \textsf{weight} carries an optional structural
coefficient used by the quantitative layer (\S\ref{sec:quant}) and ignored by the
structural queries. This encoding keeps the causal model \emph{in} the database,
co-located and jointly queryable with any other graph data, rather than in a
transient in-memory object owned by an external tool.

\subsection{The causal query catalog}\label{sec:bg:catalog}

Analysts interrogate the DAG through a small, well-established
catalog of structural queries. We group them by the reasoning they support.

\paragraph{Reachability queries} return $\mathrm{pa}$, $\mathrm{ch}$,
$\mathrm{an}$, and $\mathrm{de}$ of a variable, and test whether a directed path
$X\rightarrow\!\cdots\!\rightarrow Y$ exists. They underpin every other query.

\paragraph{Path-role queries} classify intermediate variables on the paths
between a treatment $T$ and an outcome $Y$. A \emph{mediator} lies on a directed
path $T\rightarrow\!\cdots\!\rightarrow Y$; a \emph{confounder} is a common cause,
i.e.\ a variable in $\mathrm{an}(T)\cap\mathrm{an}(Y)$; a \emph{collider} on a
path is a variable with two incoming path edges ($\rightarrow Z \leftarrow$).

\paragraph{Independence queries} test \emph{d-separation}. A path between $X$ and
$Y$ is \emph{blocked} by a conditioning set $Z$ if it contains a non-collider in
$Z$ or a collider none of whose descendants is in $Z$; $X$ and $Y$ are
d-separated given $Z$, written $X \!\perp\! Y \mid Z$, when \emph{every} path
between them is blocked~\cite{pearl2009causality}. D-separation is the graphical
criterion for conditional independence and is \emph{not} plain reachability: a
variable added to $Z$ can block one path while unblocking another through a
collider, so the test must be conditioning-aware.

\paragraph{Identification queries} decide whether the causal effect of $T$ on $Y$
can be estimated from observational data and, if so, return an \emph{adjustment
set}. A \emph{backdoor path} from $T$ to $Y$ is a path with an edge into $T$; a
set $Z$ satisfies the \emph{backdoor criterion} if it blocks every backdoor path
and contains no descendant of $T$. Enumerating minimal valid adjustment sets
issues many d-separation tests, so identification is the most traversal-intensive
query in the catalog.

Table~\ref{tab:operators} (\S\ref{sec:lang}) maps each of these to a Causal4j
language operator and the physical operator that answers it.

\subsection{Running example}\label{sec:bg:example}

We use throughout the COVID-19 causal DAG of Piccininni et
al.~\cite{piccininni2020covid}, comprising $25$ variables and $76$ directed
edges; Fig.~\ref{fig:covid} shows the fragment relevant to one query. Consider
estimating the effect of \textsf{interventions} ($T$) on
\textsf{reported\_cases} ($Y$). The causal (front-door) influence flows along
directed paths such as
$\textsf{interventions}\!\rightarrow\!\textsf{exposure\_sarscov2}\!\rightarrow\!
\textsf{onset\_covid19}\!\rightarrow\!\textsf{reported\_cases}$, whose interior
variables are mediators---and here happen to be \emph{latent}, which the engine
must nonetheless reason about. \textsf{covid19\_burden} is a common cause of both
$T$ and $Y$ (it causes \textsf{interventions} and also
\textsf{exposure\_sarscov2}), so it is a confounder that opens the backdoor path
$\textsf{interventions}\!\leftarrow\!\textsf{covid19\_burden}\!\rightarrow\!
\textsf{exposure\_sarscov2}\!\rightarrow\!\cdots\!\rightarrow\!\textsf{reported\_cases}$.
Conditioning on $\{\textsf{covid19\_burden}\}$ blocks this path and satisfies the
backdoor criterion, so it is a valid adjustment set for $T\!\rightarrow\!Y$. Every
claim in this paragraph is the output of a query in the catalog above.

\subsection{Problem statement}\label{sec:bg:problem}

We target the following data-engineering problem: \emph{support the causal query
catalog as a first-class workload inside a graph DBMS}. Concretely, a solution
must meet four requirements.
\begin{description}
  \item[R1 (Expressiveness).] A declarative interface in which each catalog query
  is a single statement, close to how analysts phrase causal questions, that
  compiles to physical operators over the graph store (\S\ref{sec:lang}).
  \item[R2 (Efficiency).] Structural queries answered without re-traversing the
  stored graph on every call, so that latency is sublinear in graph size and the
  system sustains the many independence tests an identification query issues
  (\S\ref{sec:index}).
  \item[R3 (Correctness).] Results faithful to the graphical semantics above,
  in particular a conditioning-aware d-separation test and a sound backdoor
  criterion (\S\ref{sec:index}, validated in \S\ref{sec:eval}).
  \item[R4 (Maintainability).] The accelerating index kept consistent under edge
  and variable updates without full recomputation, since curated causal graphs
  change continually (\S\ref{sec:incr}).
\end{description}
None of these is served by exporting the DAG to an external library: R2--R4 in
particular are database concerns---indexing, query processing, and incremental
maintenance---which is why we address them inside the DBMS.


%% ----------------------------------------------------------------------------
%% ============================================================================
%% Section: The Causal Query Language
%% Included from main.tex via %% ============================================================================
%% Section: The Causal Query Language
%% Included from main.tex via %% ============================================================================
%% Section: The Causal Query Language
%% Included from main.tex via \input{sections/language}
%% Requires the `causaldsl' listings style and the \TODO/\NOTE macros from main.tex.
%% Table label tab:operators is referenced from sections/background.tex (\S2.2).
%% Operator->procedure mapping is faithful to dsl/.../exec/CypherTranslator.java
%% and plugin/.../api/ProcedureNames.java.
%% ============================================================================
\section{The Causal Query Language}\label{sec:lang}

Requirement R1 asks for a declarative interface in which each catalog query is a
single statement, phrased close to how an analyst states a causal question, that
compiles to physical operators over the graph store. This section presents that
language: its surface (\S\ref{sec:lang:surface}), the mapping from each operator
to the physical operator that answers it (\S\ref{sec:lang:compile},
Table~\ref{tab:operators}), and the model registry that lets a single database
hold several causal worlds at once (\S\ref{sec:lang:registry}). The physical
operators themselves---and why they are fast---are the subject of
\S\ref{sec:index}; here we are concerned only with the interface and its
compilation.

\subsection{Design and surface syntax}\label{sec:lang:surface}

The language is a small statement-oriented DSL. A program is a sequence of
statements; each statement is one of three kinds. \emph{Definition} statements
declare variables and edges and thereby write the causal model into the graph.
\emph{Structural} statements ask the catalog queries of \S\ref{sec:bg:catalog}:
reachability, path roles, and identification. \emph{Quantitative} statements
express interventions and counterfactuals over a parameterized model
(\S\ref{sec:quant}). Two design choices follow directly from R1. First, an
operator names a causal concept, not a traversal: an analyst writes
\lstinline[style=causaldsl]|CONFOUNDERS OF T AND Y| rather than a path pattern,
so the statement reads as the question and its compilation to a graph operation
is the system's concern, not the user's. Second, every statement is
self-contained and refers to variables by their canonical
\textsf{id}, so a query is portable across any database holding a model with
those variables.

Definition statements come in a standalone and a block form. A bare
\lstinline[style=causaldsl]|DEFINE CAUSAL VARIABLE| introduces a single variable
with its observation status; a \lstinline[style=causaldsl]|DEFINE CAUSAL DAG|
block declares a model as a set of edges, optionally carrying a structural
coefficient per edge for the quantitative layer:

\begin{lstlisting}[style=causaldsl]
DEFINE CAUSAL DAG CovidModel {
  covid19_burden    -> interventions
  covid19_burden    -> exposure_sarscov2
  interventions     -> mobility
  mobility          -> exposure_sarscov2
  exposure_sarscov2 -> onset_covid19
  onset_covid19     -> reported_cases
}
\end{lstlisting}

\noindent
Against the model so defined, the structural questions of our running example
(\S\ref{sec:bg:example}) are each one statement:

\begin{lstlisting}[style=causaldsl]
CONFOUNDERS OF interventions AND reported_cases
BACKDOOR PATHS BETWEEN interventions AND reported_cases
IDENTIFY EFFECT OF interventions ON reported_cases
  RETURN ADJUSTMENT SETS EXPLAIN
\end{lstlisting}

\noindent
The first returns \textsf{covid19\_burden}; the second reports the backdoor path
it opens; the third decides identifiability and, with
\lstinline[style=causaldsl]|RETURN ADJUSTMENT SETS|, returns the valid adjustment
sets, the optional \lstinline[style=causaldsl]|EXPLAIN| attaching a
human-readable justification. The reachability and path-role operators
(\lstinline[style=causaldsl]|CAUSES OF|, \lstinline[style=causaldsl]|EFFECTS OF|,
\lstinline[style=causaldsl]|ANCESTORS OF|, \lstinline[style=causaldsl]|DESCENDANTS OF|,
\lstinline[style=causaldsl]|MEDIATORS BETWEEN|, \lstinline[style=causaldsl]|COLLIDERS BETWEEN|)
follow the same one-statement-per-question shape. The full grammar is given in
ANTLR form in the artifact; Table~\ref{tab:operators} lists the operators that
matter for what follows.

\subsection{Compilation to physical operators}\label{sec:lang:compile}

A statement is parsed to an AST and then compiled, statement by statement, to one
or more Cypher units that execute inside the DBMS. The compiler emits exactly two
kinds of unit. A \emph{definition} statement compiles to Cypher \emph{writes}
that maintain the schema of \S\ref{sec:bg:dag}: a \lstinline[style=causaldsl]|DEFINE|
becomes a \texttt{MERGE} of the \textsf{Variable} nodes and \textsf{CAUSES} edges
it introduces (a single batched \texttt{UNWIND} for the edge set), followed by a
call to \texttt{causal.compile} that refreshes the compiled snapshot
(\S\ref{sec:compile}). A \emph{structural} statement compiles to a read: a
\texttt{CALL} of the \texttt{causal.*} procedure that answers it, which consults
the snapshot rather than re-traversing the stored graph. This two-target
compilation is what keeps the model \emph{in} the database: the language never
holds the graph in client memory, and the same authoritative snapshot serves
every reader.

\begin{table}[t]
  \centering
  \caption{Language operators, their graphical semantics, and the physical
  operator each compiles to. Reachability and path-role operators map to
  \texttt{causal.*} procedures over the compiled snapshot (\S\ref{sec:index});
  d-separation is exposed as \texttt{causal.dsep} and is the inner test that
  identification issues repeatedly. Definition statements compile to Cypher
  writes plus a snapshot recompile.}
  \label{tab:operators}
  \small
  \setlength{\tabcolsep}{4pt}
  \begin{tabular}{@{}p{0.30\linewidth} p{0.34\linewidth} p{0.28\linewidth}@{}}
    \toprule
    \textbf{DSL statement} & \textbf{Graphical meaning} & \textbf{Physical operator} \\
    \midrule
    \multicolumn{3}{@{}l}{\textit{Definition}}\\
    \lstinline[style=causaldsl]|DEFINE ... DAG| & write / update model & \texttt{MERGE} + \texttt{causal.compile} \\
    \addlinespace[2pt]
    \multicolumn{3}{@{}l}{\textit{Reachability}}\\
    \lstinline[style=causaldsl]|CAUSES OF X|       & $\mathrm{pa}(X)$ & \texttt{causal.parents} \\
    \lstinline[style=causaldsl]|EFFECTS OF X|      & $\mathrm{ch}(X)$ & \texttt{causal.children} \\
    \lstinline[style=causaldsl]|ANCESTORS OF X|    & $\mathrm{an}(X)$ & \texttt{causal.ancestors} \\
    \lstinline[style=causaldsl]|DESCENDANTS OF X|  & $\mathrm{de}(X)$ & \texttt{causal.descendants} \\
    \addlinespace[2pt]
    \multicolumn{3}{@{}l}{\textit{Path roles}}\\
    \lstinline[style=causaldsl]|MEDIATORS ...|   & on a directed $T\!\to\!Y$ path & \texttt{causal.mediators} \\
    \lstinline[style=causaldsl]|COLLIDERS ...|   & $\to Z \leftarrow$ on a path & \texttt{causal.colliders} \\
    \lstinline[style=causaldsl]|CONFOUNDERS ...| & $\mathrm{an}(T)\cap\mathrm{an}(Y)$ & \texttt{causal.confounders} \\
    \lstinline[style=causaldsl]|BACKDOOR PATHS ...| & path with edge into $T$ & \texttt{causal.backdoorVariables} \\
    \addlinespace[2pt]
    \multicolumn{3}{@{}l}{\textit{Independence \& identification}}\\
    (inner test) & $X\!\perp\!Y\mid Z$ (d-sep) & \texttt{causal.dsep} \\
    \lstinline[style=causaldsl]|IDENTIFY EFFECT ...| & backdoor criterion, adj.\ sets & \texttt{causal.adjustmentCandidate} \\
    \bottomrule
  \end{tabular}
\end{table}

Structural operators map to procedures one-to-one. The reachability operators
compile to \texttt{causal.parents}, \texttt{causal.children},
\texttt{causal.ancestors}, and \texttt{causal.descendants}; the path-role
operators to \texttt{causal.mediators}, \texttt{causal.colliders}, and
\texttt{causal.confounders}; \lstinline[style=causaldsl]|BACKDOOR PATHS| to
\texttt{causal.backdoorVariables}; and \lstinline[style=causaldsl]|IDENTIFY EFFECT|
to \texttt{causal.adjustmentCandidate}, which returns whether an adjustment set
exists, the candidate sets, and---under \lstinline[style=causaldsl]|EXPLAIN|---an
explanation. D-separation is not a top-level keyword but the inner test that
identification issues repeatedly; it is exposed as the \texttt{causal.dsep}
procedure so that the cost of the traversal-intensive query (R2) is borne by a
single indexed operator (\S\ref{sec:index}) rather than by client-side path
enumeration. The quantitative statements (\lstinline[style=causaldsl]|INTERVENE|,
\lstinline[style=causaldsl]|COUNTERFACTUAL|) are handled by the quantitative
layer of \S\ref{sec:quant}; numeric intervention compiles to
\texttt{causal.linearGaussianInference}, while world construction is resolved in
the registry described next.

\subsection{Model registry and causal worlds}\label{sec:lang:registry}

A causal database is rarely a single fixed graph: a curator maintains several
models, and an intervention $do(X{=}x)$ defines a \emph{derived} model---the same
DAG with the edges into $X$ removed---that an analyst wants to name, query, and
compare against the original. The language therefore addresses models by name.
\lstinline[style=causaldsl]|DEFINE MODEL| registers a base model;
\lstinline[style=causaldsl]|USE MODEL| makes a named model (base or derived) the
active graph against which subsequent structural statements run; and
\lstinline[style=causaldsl]|RESET| returns to the base model while leaving derived
worlds registered and addressable. Context switching is resolved entirely in the
registry and emits no Cypher, so moving between worlds costs a name lookup rather
than a re-read of the graph. This registry is also the seam at which the
quantitative layer attaches: an \lstinline[style=causaldsl]|INTERVENE| statement
materializes its derived world once and registers it under a canonical name, so
that a later \lstinline[style=causaldsl]|USE MODEL| over that world reuses the
same compiled snapshot. The mechanics of snapshotting a model---base or
derived---into the layout the procedures read is the subject of the next section.

%% Requires the `causaldsl' listings style and the \TODO/\NOTE macros from main.tex.
%% Table label tab:operators is referenced from sections/background.tex (\S2.2).
%% Operator->procedure mapping is faithful to dsl/.../exec/CypherTranslator.java
%% and plugin/.../api/ProcedureNames.java.
%% ============================================================================
\section{The Causal Query Language}\label{sec:lang}

Requirement R1 asks for a declarative interface in which each catalog query is a
single statement, phrased close to how an analyst states a causal question, that
compiles to physical operators over the graph store. This section presents that
language: its surface (\S\ref{sec:lang:surface}), the mapping from each operator
to the physical operator that answers it (\S\ref{sec:lang:compile},
Table~\ref{tab:operators}), and the model registry that lets a single database
hold several causal worlds at once (\S\ref{sec:lang:registry}). The physical
operators themselves---and why they are fast---are the subject of
\S\ref{sec:index}; here we are concerned only with the interface and its
compilation.

\subsection{Design and surface syntax}\label{sec:lang:surface}

The language is a small statement-oriented DSL. A program is a sequence of
statements; each statement is one of three kinds. \emph{Definition} statements
declare variables and edges and thereby write the causal model into the graph.
\emph{Structural} statements ask the catalog queries of \S\ref{sec:bg:catalog}:
reachability, path roles, and identification. \emph{Quantitative} statements
express interventions and counterfactuals over a parameterized model
(\S\ref{sec:quant}). Two design choices follow directly from R1. First, an
operator names a causal concept, not a traversal: an analyst writes
\lstinline[style=causaldsl]|CONFOUNDERS OF T AND Y| rather than a path pattern,
so the statement reads as the question and its compilation to a graph operation
is the system's concern, not the user's. Second, every statement is
self-contained and refers to variables by their canonical
\textsf{id}, so a query is portable across any database holding a model with
those variables.

Definition statements come in a standalone and a block form. A bare
\lstinline[style=causaldsl]|DEFINE CAUSAL VARIABLE| introduces a single variable
with its observation status; a \lstinline[style=causaldsl]|DEFINE CAUSAL DAG|
block declares a model as a set of edges, optionally carrying a structural
coefficient per edge for the quantitative layer:

\begin{lstlisting}[style=causaldsl]
DEFINE CAUSAL DAG CovidModel {
  covid19_burden    -> interventions
  covid19_burden    -> exposure_sarscov2
  interventions     -> mobility
  mobility          -> exposure_sarscov2
  exposure_sarscov2 -> onset_covid19
  onset_covid19     -> reported_cases
}
\end{lstlisting}

\noindent
Against the model so defined, the structural questions of our running example
(\S\ref{sec:bg:example}) are each one statement:

\begin{lstlisting}[style=causaldsl]
CONFOUNDERS OF interventions AND reported_cases
BACKDOOR PATHS BETWEEN interventions AND reported_cases
IDENTIFY EFFECT OF interventions ON reported_cases
  RETURN ADJUSTMENT SETS EXPLAIN
\end{lstlisting}

\noindent
The first returns \textsf{covid19\_burden}; the second reports the backdoor path
it opens; the third decides identifiability and, with
\lstinline[style=causaldsl]|RETURN ADJUSTMENT SETS|, returns the valid adjustment
sets, the optional \lstinline[style=causaldsl]|EXPLAIN| attaching a
human-readable justification. The reachability and path-role operators
(\lstinline[style=causaldsl]|CAUSES OF|, \lstinline[style=causaldsl]|EFFECTS OF|,
\lstinline[style=causaldsl]|ANCESTORS OF|, \lstinline[style=causaldsl]|DESCENDANTS OF|,
\lstinline[style=causaldsl]|MEDIATORS BETWEEN|, \lstinline[style=causaldsl]|COLLIDERS BETWEEN|)
follow the same one-statement-per-question shape. The full grammar is given in
ANTLR form in the artifact; Table~\ref{tab:operators} lists the operators that
matter for what follows.

\subsection{Compilation to physical operators}\label{sec:lang:compile}

A statement is parsed to an AST and then compiled, statement by statement, to one
or more Cypher units that execute inside the DBMS. The compiler emits exactly two
kinds of unit. A \emph{definition} statement compiles to Cypher \emph{writes}
that maintain the schema of \S\ref{sec:bg:dag}: a \lstinline[style=causaldsl]|DEFINE|
becomes a \texttt{MERGE} of the \textsf{Variable} nodes and \textsf{CAUSES} edges
it introduces (a single batched \texttt{UNWIND} for the edge set), followed by a
call to \texttt{causal.compile} that refreshes the compiled snapshot
(\S\ref{sec:compile}). A \emph{structural} statement compiles to a read: a
\texttt{CALL} of the \texttt{causal.*} procedure that answers it, which consults
the snapshot rather than re-traversing the stored graph. This two-target
compilation is what keeps the model \emph{in} the database: the language never
holds the graph in client memory, and the same authoritative snapshot serves
every reader.

\begin{table}[t]
  \centering
  \caption{Language operators, their graphical semantics, and the physical
  operator each compiles to. Reachability and path-role operators map to
  \texttt{causal.*} procedures over the compiled snapshot (\S\ref{sec:index});
  d-separation is exposed as \texttt{causal.dsep} and is the inner test that
  identification issues repeatedly. Definition statements compile to Cypher
  writes plus a snapshot recompile.}
  \label{tab:operators}
  \small
  \setlength{\tabcolsep}{4pt}
  \begin{tabular}{@{}p{0.30\linewidth} p{0.34\linewidth} p{0.28\linewidth}@{}}
    \toprule
    \textbf{DSL statement} & \textbf{Graphical meaning} & \textbf{Physical operator} \\
    \midrule
    \multicolumn{3}{@{}l}{\textit{Definition}}\\
    \lstinline[style=causaldsl]|DEFINE ... DAG| & write / update model & \texttt{MERGE} + \texttt{causal.compile} \\
    \addlinespace[2pt]
    \multicolumn{3}{@{}l}{\textit{Reachability}}\\
    \lstinline[style=causaldsl]|CAUSES OF X|       & $\mathrm{pa}(X)$ & \texttt{causal.parents} \\
    \lstinline[style=causaldsl]|EFFECTS OF X|      & $\mathrm{ch}(X)$ & \texttt{causal.children} \\
    \lstinline[style=causaldsl]|ANCESTORS OF X|    & $\mathrm{an}(X)$ & \texttt{causal.ancestors} \\
    \lstinline[style=causaldsl]|DESCENDANTS OF X|  & $\mathrm{de}(X)$ & \texttt{causal.descendants} \\
    \addlinespace[2pt]
    \multicolumn{3}{@{}l}{\textit{Path roles}}\\
    \lstinline[style=causaldsl]|MEDIATORS ...|   & on a directed $T\!\to\!Y$ path & \texttt{causal.mediators} \\
    \lstinline[style=causaldsl]|COLLIDERS ...|   & $\to Z \leftarrow$ on a path & \texttt{causal.colliders} \\
    \lstinline[style=causaldsl]|CONFOUNDERS ...| & $\mathrm{an}(T)\cap\mathrm{an}(Y)$ & \texttt{causal.confounders} \\
    \lstinline[style=causaldsl]|BACKDOOR PATHS ...| & path with edge into $T$ & \texttt{causal.backdoorVariables} \\
    \addlinespace[2pt]
    \multicolumn{3}{@{}l}{\textit{Independence \& identification}}\\
    (inner test) & $X\!\perp\!Y\mid Z$ (d-sep) & \texttt{causal.dsep} \\
    \lstinline[style=causaldsl]|IDENTIFY EFFECT ...| & backdoor criterion, adj.\ sets & \texttt{causal.adjustmentCandidate} \\
    \bottomrule
  \end{tabular}
\end{table}

Structural operators map to procedures one-to-one. The reachability operators
compile to \texttt{causal.parents}, \texttt{causal.children},
\texttt{causal.ancestors}, and \texttt{causal.descendants}; the path-role
operators to \texttt{causal.mediators}, \texttt{causal.colliders}, and
\texttt{causal.confounders}; \lstinline[style=causaldsl]|BACKDOOR PATHS| to
\texttt{causal.backdoorVariables}; and \lstinline[style=causaldsl]|IDENTIFY EFFECT|
to \texttt{causal.adjustmentCandidate}, which returns whether an adjustment set
exists, the candidate sets, and---under \lstinline[style=causaldsl]|EXPLAIN|---an
explanation. D-separation is not a top-level keyword but the inner test that
identification issues repeatedly; it is exposed as the \texttt{causal.dsep}
procedure so that the cost of the traversal-intensive query (R2) is borne by a
single indexed operator (\S\ref{sec:index}) rather than by client-side path
enumeration. The quantitative statements (\lstinline[style=causaldsl]|INTERVENE|,
\lstinline[style=causaldsl]|COUNTERFACTUAL|) are handled by the quantitative
layer of \S\ref{sec:quant}; numeric intervention compiles to
\texttt{causal.linearGaussianInference}, while world construction is resolved in
the registry described next.

\subsection{Model registry and causal worlds}\label{sec:lang:registry}

A causal database is rarely a single fixed graph: a curator maintains several
models, and an intervention $do(X{=}x)$ defines a \emph{derived} model---the same
DAG with the edges into $X$ removed---that an analyst wants to name, query, and
compare against the original. The language therefore addresses models by name.
\lstinline[style=causaldsl]|DEFINE MODEL| registers a base model;
\lstinline[style=causaldsl]|USE MODEL| makes a named model (base or derived) the
active graph against which subsequent structural statements run; and
\lstinline[style=causaldsl]|RESET| returns to the base model while leaving derived
worlds registered and addressable. Context switching is resolved entirely in the
registry and emits no Cypher, so moving between worlds costs a name lookup rather
than a re-read of the graph. This registry is also the seam at which the
quantitative layer attaches: an \lstinline[style=causaldsl]|INTERVENE| statement
materializes its derived world once and registers it under a canonical name, so
that a later \lstinline[style=causaldsl]|USE MODEL| over that world reuses the
same compiled snapshot. The mechanics of snapshotting a model---base or
derived---into the layout the procedures read is the subject of the next section.

%% Requires the `causaldsl' listings style and the \TODO/\NOTE macros from main.tex.
%% Table label tab:operators is referenced from sections/background.tex (\S2.2).
%% Operator->procedure mapping is faithful to dsl/.../exec/CypherTranslator.java
%% and plugin/.../api/ProcedureNames.java.
%% ============================================================================
\section{The Causal Query Language}\label{sec:lang}

Requirement R1 asks for a declarative interface in which each catalog query is a
single statement, phrased close to how an analyst states a causal question, that
compiles to physical operators over the graph store. This section presents that
language: its surface (\S\ref{sec:lang:surface}), the mapping from each operator
to the physical operator that answers it (\S\ref{sec:lang:compile},
Table~\ref{tab:operators}), and the model registry that lets a single database
hold several causal worlds at once (\S\ref{sec:lang:registry}). The physical
operators themselves---and why they are fast---are the subject of
\S\ref{sec:index}; here we are concerned only with the interface and its
compilation.

\subsection{Design and surface syntax}\label{sec:lang:surface}

The language is a small statement-oriented DSL. A program is a sequence of
statements; each statement is one of three kinds. \emph{Definition} statements
declare variables and edges and thereby write the causal model into the graph.
\emph{Structural} statements ask the catalog queries of \S\ref{sec:bg:catalog}:
reachability, path roles, and identification. \emph{Quantitative} statements
express interventions and counterfactuals over a parameterized model
(\S\ref{sec:quant}). Two design choices follow directly from R1. First, an
operator names a causal concept, not a traversal: an analyst writes
\lstinline[style=causaldsl]|CONFOUNDERS OF T AND Y| rather than a path pattern,
so the statement reads as the question and its compilation to a graph operation
is the system's concern, not the user's. Second, every statement is
self-contained and refers to variables by their canonical
\textsf{id}, so a query is portable across any database holding a model with
those variables.

Definition statements come in a standalone and a block form. A bare
\lstinline[style=causaldsl]|DEFINE CAUSAL VARIABLE| introduces a single variable
with its observation status; a \lstinline[style=causaldsl]|DEFINE CAUSAL DAG|
block declares a model as a set of edges, optionally carrying a structural
coefficient per edge for the quantitative layer:

\begin{lstlisting}[style=causaldsl]
DEFINE CAUSAL DAG CovidModel {
  covid19_burden    -> interventions
  covid19_burden    -> exposure_sarscov2
  interventions     -> mobility
  mobility          -> exposure_sarscov2
  exposure_sarscov2 -> onset_covid19
  onset_covid19     -> reported_cases
}
\end{lstlisting}

\noindent
Against the model so defined, the structural questions of our running example
(\S\ref{sec:bg:example}) are each one statement:

\begin{lstlisting}[style=causaldsl]
CONFOUNDERS OF interventions AND reported_cases
BACKDOOR PATHS BETWEEN interventions AND reported_cases
IDENTIFY EFFECT OF interventions ON reported_cases
  RETURN ADJUSTMENT SETS EXPLAIN
\end{lstlisting}

\noindent
The first returns \textsf{covid19\_burden}; the second reports the backdoor path
it opens; the third decides identifiability and, with
\lstinline[style=causaldsl]|RETURN ADJUSTMENT SETS|, returns the valid adjustment
sets, the optional \lstinline[style=causaldsl]|EXPLAIN| attaching a
human-readable justification. The reachability and path-role operators
(\lstinline[style=causaldsl]|CAUSES OF|, \lstinline[style=causaldsl]|EFFECTS OF|,
\lstinline[style=causaldsl]|ANCESTORS OF|, \lstinline[style=causaldsl]|DESCENDANTS OF|,
\lstinline[style=causaldsl]|MEDIATORS BETWEEN|, \lstinline[style=causaldsl]|COLLIDERS BETWEEN|)
follow the same one-statement-per-question shape. The full grammar is given in
ANTLR form in the artifact; Table~\ref{tab:operators} lists the operators that
matter for what follows.

\subsection{Compilation to physical operators}\label{sec:lang:compile}

A statement is parsed to an AST and then compiled, statement by statement, to one
or more Cypher units that execute inside the DBMS. The compiler emits exactly two
kinds of unit. A \emph{definition} statement compiles to Cypher \emph{writes}
that maintain the schema of \S\ref{sec:bg:dag}: a \lstinline[style=causaldsl]|DEFINE|
becomes a \texttt{MERGE} of the \textsf{Variable} nodes and \textsf{CAUSES} edges
it introduces (a single batched \texttt{UNWIND} for the edge set), followed by a
call to \texttt{causal.compile} that refreshes the compiled snapshot
(\S\ref{sec:compile}). A \emph{structural} statement compiles to a read: a
\texttt{CALL} of the \texttt{causal.*} procedure that answers it, which consults
the snapshot rather than re-traversing the stored graph. This two-target
compilation is what keeps the model \emph{in} the database: the language never
holds the graph in client memory, and the same authoritative snapshot serves
every reader.

\begin{table}[t]
  \centering
  \caption{Language operators, their graphical semantics, and the physical
  operator each compiles to. Reachability and path-role operators map to
  \texttt{causal.*} procedures over the compiled snapshot (\S\ref{sec:index});
  d-separation is exposed as \texttt{causal.dsep} and is the inner test that
  identification issues repeatedly. Definition statements compile to Cypher
  writes plus a snapshot recompile.}
  \label{tab:operators}
  \small
  \setlength{\tabcolsep}{4pt}
  \begin{tabular}{@{}p{0.30\linewidth} p{0.34\linewidth} p{0.28\linewidth}@{}}
    \toprule
    \textbf{DSL statement} & \textbf{Graphical meaning} & \textbf{Physical operator} \\
    \midrule
    \multicolumn{3}{@{}l}{\textit{Definition}}\\
    \lstinline[style=causaldsl]|DEFINE ... DAG| & write / update model & \texttt{MERGE} + \texttt{causal.compile} \\
    \addlinespace[2pt]
    \multicolumn{3}{@{}l}{\textit{Reachability}}\\
    \lstinline[style=causaldsl]|CAUSES OF X|       & $\mathrm{pa}(X)$ & \texttt{causal.parents} \\
    \lstinline[style=causaldsl]|EFFECTS OF X|      & $\mathrm{ch}(X)$ & \texttt{causal.children} \\
    \lstinline[style=causaldsl]|ANCESTORS OF X|    & $\mathrm{an}(X)$ & \texttt{causal.ancestors} \\
    \lstinline[style=causaldsl]|DESCENDANTS OF X|  & $\mathrm{de}(X)$ & \texttt{causal.descendants} \\
    \addlinespace[2pt]
    \multicolumn{3}{@{}l}{\textit{Path roles}}\\
    \lstinline[style=causaldsl]|MEDIATORS ...|   & on a directed $T\!\to\!Y$ path & \texttt{causal.mediators} \\
    \lstinline[style=causaldsl]|COLLIDERS ...|   & $\to Z \leftarrow$ on a path & \texttt{causal.colliders} \\
    \lstinline[style=causaldsl]|CONFOUNDERS ...| & $\mathrm{an}(T)\cap\mathrm{an}(Y)$ & \texttt{causal.confounders} \\
    \lstinline[style=causaldsl]|BACKDOOR PATHS ...| & path with edge into $T$ & \texttt{causal.backdoorVariables} \\
    \addlinespace[2pt]
    \multicolumn{3}{@{}l}{\textit{Independence \& identification}}\\
    (inner test) & $X\!\perp\!Y\mid Z$ (d-sep) & \texttt{causal.dsep} \\
    \lstinline[style=causaldsl]|IDENTIFY EFFECT ...| & backdoor criterion, adj.\ sets & \texttt{causal.adjustmentCandidate} \\
    \bottomrule
  \end{tabular}
\end{table}

Structural operators map to procedures one-to-one. The reachability operators
compile to \texttt{causal.parents}, \texttt{causal.children},
\texttt{causal.ancestors}, and \texttt{causal.descendants}; the path-role
operators to \texttt{causal.mediators}, \texttt{causal.colliders}, and
\texttt{causal.confounders}; \lstinline[style=causaldsl]|BACKDOOR PATHS| to
\texttt{causal.backdoorVariables}; and \lstinline[style=causaldsl]|IDENTIFY EFFECT|
to \texttt{causal.adjustmentCandidate}, which returns whether an adjustment set
exists, the candidate sets, and---under \lstinline[style=causaldsl]|EXPLAIN|---an
explanation. D-separation is not a top-level keyword but the inner test that
identification issues repeatedly; it is exposed as the \texttt{causal.dsep}
procedure so that the cost of the traversal-intensive query (R2) is borne by a
single indexed operator (\S\ref{sec:index}) rather than by client-side path
enumeration. The quantitative statements (\lstinline[style=causaldsl]|INTERVENE|,
\lstinline[style=causaldsl]|COUNTERFACTUAL|) are handled by the quantitative
layer of \S\ref{sec:quant}; numeric intervention compiles to
\texttt{causal.linearGaussianInference}, while world construction is resolved in
the registry described next.

\subsection{Model registry and causal worlds}\label{sec:lang:registry}

A causal database is rarely a single fixed graph: a curator maintains several
models, and an intervention $do(X{=}x)$ defines a \emph{derived} model---the same
DAG with the edges into $X$ removed---that an analyst wants to name, query, and
compare against the original. The language therefore addresses models by name.
\lstinline[style=causaldsl]|DEFINE MODEL| registers a base model;
\lstinline[style=causaldsl]|USE MODEL| makes a named model (base or derived) the
active graph against which subsequent structural statements run; and
\lstinline[style=causaldsl]|RESET| returns to the base model while leaving derived
worlds registered and addressable. Context switching is resolved entirely in the
registry and emits no Cypher, so moving between worlds costs a name lookup rather
than a re-read of the graph. This registry is also the seam at which the
quantitative layer attaches: an \lstinline[style=causaldsl]|INTERVENE| statement
materializes its derived world once and registers it under a canonical name, so
that a later \lstinline[style=causaldsl]|USE MODEL| over that world reuses the
same compiled snapshot. The mechanics of snapshotting a model---base or
derived---into the layout the procedures read is the subject of the next section.


%% ----------------------------------------------------------------------------
%% ============================================================================
%% Section: Compilation -- The Causal Snapshot
%% Included from main.tex via %% ============================================================================
%% Section: Compilation -- The Causal Snapshot
%% Included from main.tex via %% ============================================================================
%% Section: Compilation -- The Causal Snapshot
%% Included from main.tex via \input{sections/compilation}
%% Faithful to GraphLoader.java, TopologicalSorter.java, CompiledCausalGraph.java,
%% RuntimeRegistry.java, DefaultCausalGraphRuntime.java.
%% ============================================================================
\section{Compilation: The Causal Snapshot}\label{sec:compile}

Structural queries are not answered against the stored property graph directly.
On \texttt{causal.compile} the graph is read once and \emph{compiled} into an
immutable in-memory snapshot that every subsequent query reads. Compilation is
the boundary between the mutable, transactional graph and the read-optimized
structure of \S\ref{sec:index}: it is where we pay the $O(V{+}E)$ scan once, so
that queries need not pay it again.

\paragraph{Loading}
The loader performs two passes inside a single read transaction. The first pass
materializes every \textsf{Variable} node with its \textsf{observed}/\textsf{latent}
flag and, when present, its linear-Gaussian parameters (intercept~$\mu_0$, noise
variance~$\sigma^2$); undeclared properties fall back to defaults ($\mu_0{=}0$,
$\sigma^2{=}1$). The second pass reads the \textsf{CAUSES} relationships with a
Cypher query that filters on \lstinline[style=causaldsl]|coalesce(r.active,true)|,
so only \emph{active} edges enter the snapshot and a curator can toggle an edge
in or out of the model without deleting it; the optional edge coefficient~$\beta$
(\textsf{weight}, default $1.0$) is read alongside for the quantitative layer.
Filtering in Cypher lets the DBMS use its own relationship access rather than
forcing a full relationship scan in Java.

\paragraph{Topological order and acyclicity}
The loader then computes a topological order of the variables with Kahn's
algorithm over the forward adjacency. This serves two purposes. It is the order
in which the index closures of \S\ref{sec:index} and the SEM recursions of
\S\ref{sec:quant} are filled, and it is the \emph{acyclicity check}: if the order
fails to cover every variable, the graph contains a directed cycle and is not a
DAG, and compilation aborts with an error rather than admitting an ill-formed
causal model. Enforcing the DAG invariant at compile time means no downstream
operator has to defend against cycles.

\paragraph{The immutable snapshot}
The result is a \emph{compiled causal graph}: variables, forward and reverse
adjacency, the topological order, and the SEM parameters, all deep-copied into
unmodifiable structures. Immutability is deliberate. Because the snapshot never
changes after construction, every reader sees a consistent view without locking,
and the accelerating index (\S\ref{sec:index}) can be built once against it and
cached rather than reconstructed per query. It is also the single source of
truth: derived facts (ancestors, adjustment sets) are exposed through
\texttt{causal.*} procedures over the snapshot rather than materialized back into
the graph as extra relationships, which would double the maintenance burden
under updates.

\paragraph{Runtime registry and consistency}
Neo4j instantiates a procedure class afresh for every call, so the snapshot
cannot live inside a procedure. It is held instead by a process-scoped registry
that owns the current snapshot in a single \texttt{volatile} reference.
Recompilation and invalidation are serialized, while all reads are lock-free:
a query captures the current reference once and works against that fixed
snapshot, so a concurrent recompile installs a new snapshot without ever exposing
a partially built one. This gives a simple, snapshot-isolated consistency model:
a query runs entirely against the model as of the compile that produced its
snapshot. When and how a new snapshot is installed---by full recompilation, or by
the incremental patch that avoids re-reading the whole graph---is the subject of
\S\ref{sec:incr}.

%% Faithful to GraphLoader.java, TopologicalSorter.java, CompiledCausalGraph.java,
%% RuntimeRegistry.java, DefaultCausalGraphRuntime.java.
%% ============================================================================
\section{Compilation: The Causal Snapshot}\label{sec:compile}

Structural queries are not answered against the stored property graph directly.
On \texttt{causal.compile} the graph is read once and \emph{compiled} into an
immutable in-memory snapshot that every subsequent query reads. Compilation is
the boundary between the mutable, transactional graph and the read-optimized
structure of \S\ref{sec:index}: it is where we pay the $O(V{+}E)$ scan once, so
that queries need not pay it again.

\paragraph{Loading}
The loader performs two passes inside a single read transaction. The first pass
materializes every \textsf{Variable} node with its \textsf{observed}/\textsf{latent}
flag and, when present, its linear-Gaussian parameters (intercept~$\mu_0$, noise
variance~$\sigma^2$); undeclared properties fall back to defaults ($\mu_0{=}0$,
$\sigma^2{=}1$). The second pass reads the \textsf{CAUSES} relationships with a
Cypher query that filters on \lstinline[style=causaldsl]|coalesce(r.active,true)|,
so only \emph{active} edges enter the snapshot and a curator can toggle an edge
in or out of the model without deleting it; the optional edge coefficient~$\beta$
(\textsf{weight}, default $1.0$) is read alongside for the quantitative layer.
Filtering in Cypher lets the DBMS use its own relationship access rather than
forcing a full relationship scan in Java.

\paragraph{Topological order and acyclicity}
The loader then computes a topological order of the variables with Kahn's
algorithm over the forward adjacency. This serves two purposes. It is the order
in which the index closures of \S\ref{sec:index} and the SEM recursions of
\S\ref{sec:quant} are filled, and it is the \emph{acyclicity check}: if the order
fails to cover every variable, the graph contains a directed cycle and is not a
DAG, and compilation aborts with an error rather than admitting an ill-formed
causal model. Enforcing the DAG invariant at compile time means no downstream
operator has to defend against cycles.

\paragraph{The immutable snapshot}
The result is a \emph{compiled causal graph}: variables, forward and reverse
adjacency, the topological order, and the SEM parameters, all deep-copied into
unmodifiable structures. Immutability is deliberate. Because the snapshot never
changes after construction, every reader sees a consistent view without locking,
and the accelerating index (\S\ref{sec:index}) can be built once against it and
cached rather than reconstructed per query. It is also the single source of
truth: derived facts (ancestors, adjustment sets) are exposed through
\texttt{causal.*} procedures over the snapshot rather than materialized back into
the graph as extra relationships, which would double the maintenance burden
under updates.

\paragraph{Runtime registry and consistency}
Neo4j instantiates a procedure class afresh for every call, so the snapshot
cannot live inside a procedure. It is held instead by a process-scoped registry
that owns the current snapshot in a single \texttt{volatile} reference.
Recompilation and invalidation are serialized, while all reads are lock-free:
a query captures the current reference once and works against that fixed
snapshot, so a concurrent recompile installs a new snapshot without ever exposing
a partially built one. This gives a simple, snapshot-isolated consistency model:
a query runs entirely against the model as of the compile that produced its
snapshot. When and how a new snapshot is installed---by full recompilation, or by
the incremental patch that avoids re-reading the whole graph---is the subject of
\S\ref{sec:incr}.

%% Faithful to GraphLoader.java, TopologicalSorter.java, CompiledCausalGraph.java,
%% RuntimeRegistry.java, DefaultCausalGraphRuntime.java.
%% ============================================================================
\section{Compilation: The Causal Snapshot}\label{sec:compile}

Structural queries are not answered against the stored property graph directly.
On \texttt{causal.compile} the graph is read once and \emph{compiled} into an
immutable in-memory snapshot that every subsequent query reads. Compilation is
the boundary between the mutable, transactional graph and the read-optimized
structure of \S\ref{sec:index}: it is where we pay the $O(V{+}E)$ scan once, so
that queries need not pay it again.

\paragraph{Loading}
The loader performs two passes inside a single read transaction. The first pass
materializes every \textsf{Variable} node with its \textsf{observed}/\textsf{latent}
flag and, when present, its linear-Gaussian parameters (intercept~$\mu_0$, noise
variance~$\sigma^2$); undeclared properties fall back to defaults ($\mu_0{=}0$,
$\sigma^2{=}1$). The second pass reads the \textsf{CAUSES} relationships with a
Cypher query that filters on \lstinline[style=causaldsl]|coalesce(r.active,true)|,
so only \emph{active} edges enter the snapshot and a curator can toggle an edge
in or out of the model without deleting it; the optional edge coefficient~$\beta$
(\textsf{weight}, default $1.0$) is read alongside for the quantitative layer.
Filtering in Cypher lets the DBMS use its own relationship access rather than
forcing a full relationship scan in Java.

\paragraph{Topological order and acyclicity}
The loader then computes a topological order of the variables with Kahn's
algorithm over the forward adjacency. This serves two purposes. It is the order
in which the index closures of \S\ref{sec:index} and the SEM recursions of
\S\ref{sec:quant} are filled, and it is the \emph{acyclicity check}: if the order
fails to cover every variable, the graph contains a directed cycle and is not a
DAG, and compilation aborts with an error rather than admitting an ill-formed
causal model. Enforcing the DAG invariant at compile time means no downstream
operator has to defend against cycles.

\paragraph{The immutable snapshot}
The result is a \emph{compiled causal graph}: variables, forward and reverse
adjacency, the topological order, and the SEM parameters, all deep-copied into
unmodifiable structures. Immutability is deliberate. Because the snapshot never
changes after construction, every reader sees a consistent view without locking,
and the accelerating index (\S\ref{sec:index}) can be built once against it and
cached rather than reconstructed per query. It is also the single source of
truth: derived facts (ancestors, adjustment sets) are exposed through
\texttt{causal.*} procedures over the snapshot rather than materialized back into
the graph as extra relationships, which would double the maintenance burden
under updates.

\paragraph{Runtime registry and consistency}
Neo4j instantiates a procedure class afresh for every call, so the snapshot
cannot live inside a procedure. It is held instead by a process-scoped registry
that owns the current snapshot in a single \texttt{volatile} reference.
Recompilation and invalidation are serialized, while all reads are lock-free:
a query captures the current reference once and works against that fixed
snapshot, so a concurrent recompile installs a new snapshot without ever exposing
a partially built one. This gives a simple, snapshot-isolated consistency model:
a query runs entirely against the model as of the compile that produced its
snapshot. When and how a new snapshot is installed---by full recompilation, or by
the incremental patch that avoids re-reading the whole graph---is the subject of
\S\ref{sec:incr}.


%% ----------------------------------------------------------------------------
%% ============================================================================
%% Section: Indexing and Query Processing  (CORE)
%% Included from main.tex via %% ============================================================================
%% Section: Indexing and Query Processing  (CORE)
%% Included from main.tex via %% ============================================================================
%% Section: Indexing and Query Processing  (CORE)
%% Included from main.tex via \input{sections/indexing}
%% Owns Table tab:complexity and Fig fig:index.
%% Design faithful to the Tier-2 layout planned in CompiledCausalGraph.java and
%% the algorithms in DefaultCausalGraphRuntime.java / PathAnalyzer.java /
%% ReachabilityIndex.java. See the status note (\NOTE) on implementation state.
%% ============================================================================
\section{Indexing and Query Processing}\label{sec:index}

This section is the core of the paper. It presents the physical layout the causal
snapshot compiles to (\S\ref{sec:index:layout}), the precomputed reachability
bitsets that accelerate every structural query (\S\ref{sec:index:bitsets}), and
the query-processing algorithms built on them: reachability and path roles
(\S\ref{sec:index:reach}), conditioning-aware d-separation
(\S\ref{sec:index:dsep}), and adjustment-set enumeration
(\S\ref{sec:index:adj}). Table~\ref{tab:complexity} summarizes the per-operator
cost against a map-of-sets baseline.

\NOTE{Implementation status (remove before CR; kept here per the honest-gaps
brief). The algorithms in \S\ref{sec:index:reach}--\S\ref{sec:index:adj} are
implemented and unit-tested (Bayes-Ball d-separation, backdoor adjustment
enumeration, path roles). The \emph{dense CSR + bitset} physical layout of
\S\ref{sec:index:layout}--\S\ref{sec:index:bitsets} is specified but not yet the
running representation: the current prototype answers these queries over a
map-of-sets adjacency with DFS closures (the $O(V{+}E)$ baseline column of
Table~\ref{tab:complexity}). The bitset layout is what the \S\ref{sec:eval}
speedups will measure; the ``Index query'' column states the design target, not a
measured result. Do not claim the speedup as achieved until the eval lands.}

\subsection{Physical layout}\label{sec:index:layout}

The snapshot's variables are assigned dense integer indices $0..n{-}1$ via a
\textsf{name}$\to$\textsf{index} map, so every downstream structure is an array
rather than a hash map (Fig.~\ref{fig:index}). Edges are stored as compressed
adjacency: a forward and a reverse \emph{CSR}-style array pair
($\textsf{forwardAdj}[i]$ lists the child indices of $i$, $\textsf{reverseAdj}[i]$
its parents), giving cache-friendly, allocation-free neighbor iteration. The
topological order is stored as an \texttt{int[]} $\textsf{topoOrder}$, and each
node carries its \emph{layer} $\textsf{topoLayer}[i]$ (0 for sources), which lets
a traversal prune any neighbor that cannot lie on a path between two nodes by a
simple layer comparison. This replaces the baseline representation---a
\lstinline[style=causaldsl]|Map<String,Set<String>>| adjacency with
\texttt{String}-keyed lookups and per-query DFS---whose every structural query
costs an $O(V{+}E)$ traversal with hashing overhead.

\subsection{Reachability bitsets}\label{sec:index:bitsets}

On top of the dense layout we precompute, for every variable, the set of its
ancestors and the set of its descendants as fixed-width bitsets
$\textsf{ancestorBits}[i],\textsf{descendantBits}[i]\in\{0,1\}^{n}$, packed into
$\lceil n/64\rceil$ machine words. The ancestor closures are filled in a single
reverse-topological sweep using the recurrence
\[
  \textsf{ancestorBits}[i] \;=\; \bigcup_{p\,\in\,\mathrm{pa}(i)}
      \bigl(\textsf{ancestorBits}[p] \;\cup\; \{p\}\bigr),
\]
and descendants symmetrically in forward-topological order; each is one
$O(V{+}E)$ pass at compile time whose cost is amortized over all later queries.
The closures occupy $O(V^2/64)$ words---the deliberate space-for-time trade that
turns reachability from traversal into a word-parallel bit test. Sparse or very
large DAGs can instead keep only the CSR layout and answer reachability by a
layer-pruned DFS; \S\ref{sec:eval} characterizes the memory/latency crossover
between the two.

\begin{figure}[t]\centering
  \resizebox{\linewidth}{!}{\input{figures/index-layout}}
  \caption{Compiled index layout for a small fragment. Variables map to dense
  indices; forward/reverse CSR arrays store adjacency; a topological order and
  per-node layer drive pruning; and per-node ancestor/descendant bitsets
  ($\lceil n/64\rceil$ words each) are filled in reverse-/forward-topological
  order. Structural queries become word-parallel bit operations over the
  bitsets.}
  \label{fig:index}
\end{figure}

\subsection{Reachability and path-role queries}\label{sec:index:reach}

With the closures in hand the reachability operators are immediate:
$\mathrm{an}(X)$ and $\mathrm{de}(X)$ are a single bitset read, and
$\textsf{hasDirectedPath}(X,Y)$ is one bit test---is $Y$ set in
$\textsf{descendantBits}[X]$---in $O(1)$. The path-role operators reduce to
bitset algebra over these closures. Confounders of $T$ and $Y$ are the common
ancestors $\textsf{ancestorBits}[T]\cap\textsf{ancestorBits}[Y]$; mediators of
$T\!\to\!Y$ are $\textsf{descendantBits}[T]\cap\textsf{ancestorBits}[Y]$; and
colliders are the nodes whose parents include one reachable from $X$ and one from
$Y$ while the node itself ancestors neither---again a handful of intersections
per candidate. Each is $O(V/64)$ word operations rather than the baseline's
repeated DFS.

\subsection{Conditioning-aware d-separation}\label{sec:index:dsep}

D-separation is \emph{not} reachability: adding a variable to the conditioning set
$Z$ can block one path and, if that variable is a collider or an ancestor of one,
unblock another. The test must therefore follow edge orientation and the role of
each node under $Z$, which we do with the Bayes-Ball procedure. A search
propagates ``balls'' from $X$, tracking for each visited node whether it was
entered from a parent or from a child; at each node the admissible continuations
depend on whether the node is in $Z$: an observed non-collider blocks a ball
arriving along the chain/fork, while a collider (or a node with an observed
descendant) passes a ball arriving from two children. $X$ and $Y$ are
d-separated given $Z$ exactly when no ball reaches $Y$. The procedure is
conditioning-aware by construction and, with conditioning and frontier sets
represented as bitsets, its per-step membership tests are $O(1)$ and the whole
test $O(V/64)$ in the frontier. A single flag---cutting the outgoing edges of the
treatment---turns the same routine into the backdoor test used next.

\subsection{Adjustment-set enumeration}\label{sec:index:adj}

Identification asks whether the effect of $T$ on $Y$ is estimable by covariate
adjustment and, if so, for a valid set. A set $Z$ satisfies the \emph{backdoor
criterion} when it contains no descendant of $T$ and blocks every backdoor path
from $T$ to $Y$; the second condition is exactly the d-separation test of
\S\ref{sec:index:dsep} with $T$'s outgoing edges cut. Enumeration first restricts
the candidate pool to \emph{eligible} variables---those that are not $T$, not $Y$,
and (via $\textsf{descendantBits}[T]$) not descendants of $T$, so the
no-descendant clause holds for free. It then searches candidate sets in
increasing size and returns the first non-empty level of blocking sets, yielding
\emph{minimal} adjustment sets; the empty set is tested first, so a query with no
open backdoor path returns immediately. This is the most traversal-intensive
operator in the catalog---it issues one backdoor (d-separation) test per
candidate set---which is precisely why casting the inner test as an $O(V/64)$
bitset operation, rather than a graph traversal, is what makes identification
scale (R2).

\begin{table}[t]\centering
\caption{Per-operator cost: map-of-sets $+$ DFS baseline vs.\ the bitset index.
$k$ is the number of candidate sets an identification query tests.}
\label{tab:complexity}
\begin{tabular}{@{}lccc@{}}
\toprule
Operator & Baseline query & Index query & Index space \\
\midrule
ancestors / descendants & $O(V{+}E)$ & $O(V/64)$ & $O(V^2/64)$ \\
hasDirectedPath         & $O(V{+}E)$ & $O(1)$     & --- \\
mediators/conf./coll.   & $O(V{+}E)$ & $O(V/64)$  & --- \\
d-separation            & $O(V{+}E)$ & $O(V/64)$  & --- \\
adjustment sets         & $O(k(V{+}E))$ & $O(k\,V/64)$ & --- \\
\bottomrule
\end{tabular}
\end{table}

%% Owns Table tab:complexity and Fig fig:index.
%% Design faithful to the Tier-2 layout planned in CompiledCausalGraph.java and
%% the algorithms in DefaultCausalGraphRuntime.java / PathAnalyzer.java /
%% ReachabilityIndex.java. See the status note (\NOTE) on implementation state.
%% ============================================================================
\section{Indexing and Query Processing}\label{sec:index}

This section is the core of the paper. It presents the physical layout the causal
snapshot compiles to (\S\ref{sec:index:layout}), the precomputed reachability
bitsets that accelerate every structural query (\S\ref{sec:index:bitsets}), and
the query-processing algorithms built on them: reachability and path roles
(\S\ref{sec:index:reach}), conditioning-aware d-separation
(\S\ref{sec:index:dsep}), and adjustment-set enumeration
(\S\ref{sec:index:adj}). Table~\ref{tab:complexity} summarizes the per-operator
cost against a map-of-sets baseline.

\NOTE{Implementation status (remove before CR; kept here per the honest-gaps
brief). The algorithms in \S\ref{sec:index:reach}--\S\ref{sec:index:adj} are
implemented and unit-tested (Bayes-Ball d-separation, backdoor adjustment
enumeration, path roles). The \emph{dense CSR + bitset} physical layout of
\S\ref{sec:index:layout}--\S\ref{sec:index:bitsets} is specified but not yet the
running representation: the current prototype answers these queries over a
map-of-sets adjacency with DFS closures (the $O(V{+}E)$ baseline column of
Table~\ref{tab:complexity}). The bitset layout is what the \S\ref{sec:eval}
speedups will measure; the ``Index query'' column states the design target, not a
measured result. Do not claim the speedup as achieved until the eval lands.}

\subsection{Physical layout}\label{sec:index:layout}

The snapshot's variables are assigned dense integer indices $0..n{-}1$ via a
\textsf{name}$\to$\textsf{index} map, so every downstream structure is an array
rather than a hash map (Fig.~\ref{fig:index}). Edges are stored as compressed
adjacency: a forward and a reverse \emph{CSR}-style array pair
($\textsf{forwardAdj}[i]$ lists the child indices of $i$, $\textsf{reverseAdj}[i]$
its parents), giving cache-friendly, allocation-free neighbor iteration. The
topological order is stored as an \texttt{int[]} $\textsf{topoOrder}$, and each
node carries its \emph{layer} $\textsf{topoLayer}[i]$ (0 for sources), which lets
a traversal prune any neighbor that cannot lie on a path between two nodes by a
simple layer comparison. This replaces the baseline representation---a
\lstinline[style=causaldsl]|Map<String,Set<String>>| adjacency with
\texttt{String}-keyed lookups and per-query DFS---whose every structural query
costs an $O(V{+}E)$ traversal with hashing overhead.

\subsection{Reachability bitsets}\label{sec:index:bitsets}

On top of the dense layout we precompute, for every variable, the set of its
ancestors and the set of its descendants as fixed-width bitsets
$\textsf{ancestorBits}[i],\textsf{descendantBits}[i]\in\{0,1\}^{n}$, packed into
$\lceil n/64\rceil$ machine words. The ancestor closures are filled in a single
reverse-topological sweep using the recurrence
\[
  \textsf{ancestorBits}[i] \;=\; \bigcup_{p\,\in\,\mathrm{pa}(i)}
      \bigl(\textsf{ancestorBits}[p] \;\cup\; \{p\}\bigr),
\]
and descendants symmetrically in forward-topological order; each is one
$O(V{+}E)$ pass at compile time whose cost is amortized over all later queries.
The closures occupy $O(V^2/64)$ words---the deliberate space-for-time trade that
turns reachability from traversal into a word-parallel bit test. Sparse or very
large DAGs can instead keep only the CSR layout and answer reachability by a
layer-pruned DFS; \S\ref{sec:eval} characterizes the memory/latency crossover
between the two.

\begin{figure}[t]\centering
  \resizebox{\linewidth}{!}{\input{figures/index-layout}}
  \caption{Compiled index layout for a small fragment. Variables map to dense
  indices; forward/reverse CSR arrays store adjacency; a topological order and
  per-node layer drive pruning; and per-node ancestor/descendant bitsets
  ($\lceil n/64\rceil$ words each) are filled in reverse-/forward-topological
  order. Structural queries become word-parallel bit operations over the
  bitsets.}
  \label{fig:index}
\end{figure}

\subsection{Reachability and path-role queries}\label{sec:index:reach}

With the closures in hand the reachability operators are immediate:
$\mathrm{an}(X)$ and $\mathrm{de}(X)$ are a single bitset read, and
$\textsf{hasDirectedPath}(X,Y)$ is one bit test---is $Y$ set in
$\textsf{descendantBits}[X]$---in $O(1)$. The path-role operators reduce to
bitset algebra over these closures. Confounders of $T$ and $Y$ are the common
ancestors $\textsf{ancestorBits}[T]\cap\textsf{ancestorBits}[Y]$; mediators of
$T\!\to\!Y$ are $\textsf{descendantBits}[T]\cap\textsf{ancestorBits}[Y]$; and
colliders are the nodes whose parents include one reachable from $X$ and one from
$Y$ while the node itself ancestors neither---again a handful of intersections
per candidate. Each is $O(V/64)$ word operations rather than the baseline's
repeated DFS.

\subsection{Conditioning-aware d-separation}\label{sec:index:dsep}

D-separation is \emph{not} reachability: adding a variable to the conditioning set
$Z$ can block one path and, if that variable is a collider or an ancestor of one,
unblock another. The test must therefore follow edge orientation and the role of
each node under $Z$, which we do with the Bayes-Ball procedure. A search
propagates ``balls'' from $X$, tracking for each visited node whether it was
entered from a parent or from a child; at each node the admissible continuations
depend on whether the node is in $Z$: an observed non-collider blocks a ball
arriving along the chain/fork, while a collider (or a node with an observed
descendant) passes a ball arriving from two children. $X$ and $Y$ are
d-separated given $Z$ exactly when no ball reaches $Y$. The procedure is
conditioning-aware by construction and, with conditioning and frontier sets
represented as bitsets, its per-step membership tests are $O(1)$ and the whole
test $O(V/64)$ in the frontier. A single flag---cutting the outgoing edges of the
treatment---turns the same routine into the backdoor test used next.

\subsection{Adjustment-set enumeration}\label{sec:index:adj}

Identification asks whether the effect of $T$ on $Y$ is estimable by covariate
adjustment and, if so, for a valid set. A set $Z$ satisfies the \emph{backdoor
criterion} when it contains no descendant of $T$ and blocks every backdoor path
from $T$ to $Y$; the second condition is exactly the d-separation test of
\S\ref{sec:index:dsep} with $T$'s outgoing edges cut. Enumeration first restricts
the candidate pool to \emph{eligible} variables---those that are not $T$, not $Y$,
and (via $\textsf{descendantBits}[T]$) not descendants of $T$, so the
no-descendant clause holds for free. It then searches candidate sets in
increasing size and returns the first non-empty level of blocking sets, yielding
\emph{minimal} adjustment sets; the empty set is tested first, so a query with no
open backdoor path returns immediately. This is the most traversal-intensive
operator in the catalog---it issues one backdoor (d-separation) test per
candidate set---which is precisely why casting the inner test as an $O(V/64)$
bitset operation, rather than a graph traversal, is what makes identification
scale (R2).

\begin{table}[t]\centering
\caption{Per-operator cost: map-of-sets $+$ DFS baseline vs.\ the bitset index.
$k$ is the number of candidate sets an identification query tests.}
\label{tab:complexity}
\begin{tabular}{@{}lccc@{}}
\toprule
Operator & Baseline query & Index query & Index space \\
\midrule
ancestors / descendants & $O(V{+}E)$ & $O(V/64)$ & $O(V^2/64)$ \\
hasDirectedPath         & $O(V{+}E)$ & $O(1)$     & --- \\
mediators/conf./coll.   & $O(V{+}E)$ & $O(V/64)$  & --- \\
d-separation            & $O(V{+}E)$ & $O(V/64)$  & --- \\
adjustment sets         & $O(k(V{+}E))$ & $O(k\,V/64)$ & --- \\
\bottomrule
\end{tabular}
\end{table}

%% Owns Table tab:complexity and Fig fig:index.
%% Design faithful to the Tier-2 layout planned in CompiledCausalGraph.java and
%% the algorithms in DefaultCausalGraphRuntime.java / PathAnalyzer.java /
%% ReachabilityIndex.java. See the status note (\NOTE) on implementation state.
%% ============================================================================
\section{Indexing and Query Processing}\label{sec:index}

This section is the core of the paper. It presents the physical layout the causal
snapshot compiles to (\S\ref{sec:index:layout}), the precomputed reachability
bitsets that accelerate every structural query (\S\ref{sec:index:bitsets}), and
the query-processing algorithms built on them: reachability and path roles
(\S\ref{sec:index:reach}), conditioning-aware d-separation
(\S\ref{sec:index:dsep}), and adjustment-set enumeration
(\S\ref{sec:index:adj}). Table~\ref{tab:complexity} summarizes the per-operator
cost against a map-of-sets baseline.

\NOTE{Implementation status (remove before CR; kept here per the honest-gaps
brief). The algorithms in \S\ref{sec:index:reach}--\S\ref{sec:index:adj} are
implemented and unit-tested (Bayes-Ball d-separation, backdoor adjustment
enumeration, path roles). The \emph{dense CSR + bitset} physical layout of
\S\ref{sec:index:layout}--\S\ref{sec:index:bitsets} is specified but not yet the
running representation: the current prototype answers these queries over a
map-of-sets adjacency with DFS closures (the $O(V{+}E)$ baseline column of
Table~\ref{tab:complexity}). The bitset layout is what the \S\ref{sec:eval}
speedups will measure; the ``Index query'' column states the design target, not a
measured result. Do not claim the speedup as achieved until the eval lands.}

\subsection{Physical layout}\label{sec:index:layout}

The snapshot's variables are assigned dense integer indices $0..n{-}1$ via a
\textsf{name}$\to$\textsf{index} map, so every downstream structure is an array
rather than a hash map (Fig.~\ref{fig:index}). Edges are stored as compressed
adjacency: a forward and a reverse \emph{CSR}-style array pair
($\textsf{forwardAdj}[i]$ lists the child indices of $i$, $\textsf{reverseAdj}[i]$
its parents), giving cache-friendly, allocation-free neighbor iteration. The
topological order is stored as an \texttt{int[]} $\textsf{topoOrder}$, and each
node carries its \emph{layer} $\textsf{topoLayer}[i]$ (0 for sources), which lets
a traversal prune any neighbor that cannot lie on a path between two nodes by a
simple layer comparison. This replaces the baseline representation---a
\lstinline[style=causaldsl]|Map<String,Set<String>>| adjacency with
\texttt{String}-keyed lookups and per-query DFS---whose every structural query
costs an $O(V{+}E)$ traversal with hashing overhead.

\subsection{Reachability bitsets}\label{sec:index:bitsets}

On top of the dense layout we precompute, for every variable, the set of its
ancestors and the set of its descendants as fixed-width bitsets
$\textsf{ancestorBits}[i],\textsf{descendantBits}[i]\in\{0,1\}^{n}$, packed into
$\lceil n/64\rceil$ machine words. The ancestor closures are filled in a single
reverse-topological sweep using the recurrence
\[
  \textsf{ancestorBits}[i] \;=\; \bigcup_{p\,\in\,\mathrm{pa}(i)}
      \bigl(\textsf{ancestorBits}[p] \;\cup\; \{p\}\bigr),
\]
and descendants symmetrically in forward-topological order; each is one
$O(V{+}E)$ pass at compile time whose cost is amortized over all later queries.
The closures occupy $O(V^2/64)$ words---the deliberate space-for-time trade that
turns reachability from traversal into a word-parallel bit test. Sparse or very
large DAGs can instead keep only the CSR layout and answer reachability by a
layer-pruned DFS; \S\ref{sec:eval} characterizes the memory/latency crossover
between the two.

\begin{figure}[t]\centering
  \resizebox{\linewidth}{!}{\input{figures/index-layout}}
  \caption{Compiled index layout for a small fragment. Variables map to dense
  indices; forward/reverse CSR arrays store adjacency; a topological order and
  per-node layer drive pruning; and per-node ancestor/descendant bitsets
  ($\lceil n/64\rceil$ words each) are filled in reverse-/forward-topological
  order. Structural queries become word-parallel bit operations over the
  bitsets.}
  \label{fig:index}
\end{figure}

\subsection{Reachability and path-role queries}\label{sec:index:reach}

With the closures in hand the reachability operators are immediate:
$\mathrm{an}(X)$ and $\mathrm{de}(X)$ are a single bitset read, and
$\textsf{hasDirectedPath}(X,Y)$ is one bit test---is $Y$ set in
$\textsf{descendantBits}[X]$---in $O(1)$. The path-role operators reduce to
bitset algebra over these closures. Confounders of $T$ and $Y$ are the common
ancestors $\textsf{ancestorBits}[T]\cap\textsf{ancestorBits}[Y]$; mediators of
$T\!\to\!Y$ are $\textsf{descendantBits}[T]\cap\textsf{ancestorBits}[Y]$; and
colliders are the nodes whose parents include one reachable from $X$ and one from
$Y$ while the node itself ancestors neither---again a handful of intersections
per candidate. Each is $O(V/64)$ word operations rather than the baseline's
repeated DFS.

\subsection{Conditioning-aware d-separation}\label{sec:index:dsep}

D-separation is \emph{not} reachability: adding a variable to the conditioning set
$Z$ can block one path and, if that variable is a collider or an ancestor of one,
unblock another. The test must therefore follow edge orientation and the role of
each node under $Z$, which we do with the Bayes-Ball procedure. A search
propagates ``balls'' from $X$, tracking for each visited node whether it was
entered from a parent or from a child; at each node the admissible continuations
depend on whether the node is in $Z$: an observed non-collider blocks a ball
arriving along the chain/fork, while a collider (or a node with an observed
descendant) passes a ball arriving from two children. $X$ and $Y$ are
d-separated given $Z$ exactly when no ball reaches $Y$. The procedure is
conditioning-aware by construction and, with conditioning and frontier sets
represented as bitsets, its per-step membership tests are $O(1)$ and the whole
test $O(V/64)$ in the frontier. A single flag---cutting the outgoing edges of the
treatment---turns the same routine into the backdoor test used next.

\subsection{Adjustment-set enumeration}\label{sec:index:adj}

Identification asks whether the effect of $T$ on $Y$ is estimable by covariate
adjustment and, if so, for a valid set. A set $Z$ satisfies the \emph{backdoor
criterion} when it contains no descendant of $T$ and blocks every backdoor path
from $T$ to $Y$; the second condition is exactly the d-separation test of
\S\ref{sec:index:dsep} with $T$'s outgoing edges cut. Enumeration first restricts
the candidate pool to \emph{eligible} variables---those that are not $T$, not $Y$,
and (via $\textsf{descendantBits}[T]$) not descendants of $T$, so the
no-descendant clause holds for free. It then searches candidate sets in
increasing size and returns the first non-empty level of blocking sets, yielding
\emph{minimal} adjustment sets; the empty set is tested first, so a query with no
open backdoor path returns immediately. This is the most traversal-intensive
operator in the catalog---it issues one backdoor (d-separation) test per
candidate set---which is precisely why casting the inner test as an $O(V/64)$
bitset operation, rather than a graph traversal, is what makes identification
scale (R2).

\begin{table}[t]\centering
\caption{Per-operator cost: map-of-sets $+$ DFS baseline vs.\ the bitset index.
$k$ is the number of candidate sets an identification query tests.}
\label{tab:complexity}
\begin{tabular}{@{}lccc@{}}
\toprule
Operator & Baseline query & Index query & Index space \\
\midrule
ancestors / descendants & $O(V{+}E)$ & $O(V/64)$ & $O(V^2/64)$ \\
hasDirectedPath         & $O(V{+}E)$ & $O(1)$     & --- \\
mediators/conf./coll.   & $O(V{+}E)$ & $O(V/64)$  & --- \\
d-separation            & $O(V{+}E)$ & $O(V/64)$  & --- \\
adjustment sets         & $O(k(V{+}E))$ & $O(k\,V/64)$ & --- \\
\bottomrule
\end{tabular}
\end{table}


%% ----------------------------------------------------------------------------
%% ============================================================================
%% Section: Incremental Index Maintenance
%% Included from main.tex via %% ============================================================================
%% Section: Incremental Index Maintenance
%% Included from main.tex via %% ============================================================================
%% Section: Incremental Index Maintenance
%% Included from main.tex via \input{sections/incremental}
%% Faithful to lifecycle/CausalGraphChangeListener.java, ChangeBatch.java,
%% CausalGraphLifecycle.java. See the status note on what is implemented vs.
%% designed.
%% ============================================================================
\section{Incremental Index Maintenance}\label{sec:incr}

A curated causal graph is edited continually---an edge is added as evidence
accrues, a variable is reclassified as latent, an edge is deactivated pending
review. Requirement R4 asks that the compiled index survive these edits without
a full recompilation after every change. This section describes how Causal4j
observes graph mutations and keeps the snapshot consistent.

\paragraph{Observing changes}
The plugin registers a transaction-event listener, scoped to the database for its
lifetime, that sees every committing write. Before a transaction commits, the
listener inspects the transaction data and, if it touches the causal schema,
distills the change into a compact \emph{change batch}: the ids of created and
deleted \textsf{Variable} nodes and the endpoint pairs of created and deleted
\textsf{CAUSES} edges. Transactions that touch no causal schema are rejected by a
cheap guard and cost nothing. The pre-commit hook is also the natural place to
\emph{validate}: because a well-formed causal model must remain acyclic, an edge
insertion that would close a cycle can be rejected here, aborting the offending
transaction rather than admitting an invalid model into the index. The batch is
handed to the post-commit hook, which applies it to the runtime once the write is
durable.

\paragraph{Incremental patch}
Given a change batch, the closures of \S\ref{sec:index:bitsets} can be maintained
locally rather than rebuilt. Inserting an edge $u\!\to\!v$ (having checked it
introduces no cycle) extends the CSR adjacency and then propagates closures over
only the \emph{affected cone}: $u$ and its ancestors gain $v$ and its descendants
as new descendants, and symmetrically $v$ and its descendants gain $u$ and its
ancestors as new ancestors. Concretely, $\textsf{descendantBits}[u]\cup\{v\}$ is
OR-ed into the descendant bitset of every ancestor of $u$, and the dual update is
applied for ancestors; nodes outside the cone are untouched, so the work is
bounded by the size of the affected region rather than by $|V|$. Adding or
removing a variable allocates or frees a dense index and scrubs its incident
edges. A property flip on an edge's \textsf{active} flag is treated as an edge
insertion or deletion, and a flip of a variable's \textsf{latent} flag is a
property-only update that leaves the topology, and therefore the closures,
untouched.

\paragraph{When to fall back}
Not every edit patches cheaply. Edge \emph{deletion} can shrink ancestor and
descendant closures in non-local ways---a single removed edge may sever
reachability for many pairs---so when a deletion invalidates cached closures the
sound choice is to recompute the affected cone or, in the worst case, fall back to
a full recompilation. The system therefore treats incremental maintenance as an
optimization over the always-correct recompile: additive changes and property
flips patch in place, while topology-shrinking changes fall back when a bounded
patch would be unsound. \S\ref{sec:eval} measures the crossover---batch size and
edit type at which incremental maintenance beats recompilation
(Fig.~\ref{fig:incr}).

\NOTE{Implementation status (remove before CR; honest-gaps brief). The
transaction-listener plumbing---schema-change detection, the change-batch diff of
added/removed variables and edges, database-scoped registration/teardown---is
implemented. The post-commit hook currently applies the always-correct
\emph{fallback}: it invalidates the snapshot so the next query recompiles. The
bounded in-place patch described above (the OR-in over the affected cone) and the
pre-commit cycle/active/latent validators are specified but not yet wired; the
change batch is deliberately the context object through which they will be
applied. Frame R4 as ``designed, plumbing in place, patch pending'' until the
applyChange path and its eval land.}

%% Faithful to lifecycle/CausalGraphChangeListener.java, ChangeBatch.java,
%% CausalGraphLifecycle.java. See the status note on what is implemented vs.
%% designed.
%% ============================================================================
\section{Incremental Index Maintenance}\label{sec:incr}

A curated causal graph is edited continually---an edge is added as evidence
accrues, a variable is reclassified as latent, an edge is deactivated pending
review. Requirement R4 asks that the compiled index survive these edits without
a full recompilation after every change. This section describes how Causal4j
observes graph mutations and keeps the snapshot consistent.

\paragraph{Observing changes}
The plugin registers a transaction-event listener, scoped to the database for its
lifetime, that sees every committing write. Before a transaction commits, the
listener inspects the transaction data and, if it touches the causal schema,
distills the change into a compact \emph{change batch}: the ids of created and
deleted \textsf{Variable} nodes and the endpoint pairs of created and deleted
\textsf{CAUSES} edges. Transactions that touch no causal schema are rejected by a
cheap guard and cost nothing. The pre-commit hook is also the natural place to
\emph{validate}: because a well-formed causal model must remain acyclic, an edge
insertion that would close a cycle can be rejected here, aborting the offending
transaction rather than admitting an invalid model into the index. The batch is
handed to the post-commit hook, which applies it to the runtime once the write is
durable.

\paragraph{Incremental patch}
Given a change batch, the closures of \S\ref{sec:index:bitsets} can be maintained
locally rather than rebuilt. Inserting an edge $u\!\to\!v$ (having checked it
introduces no cycle) extends the CSR adjacency and then propagates closures over
only the \emph{affected cone}: $u$ and its ancestors gain $v$ and its descendants
as new descendants, and symmetrically $v$ and its descendants gain $u$ and its
ancestors as new ancestors. Concretely, $\textsf{descendantBits}[u]\cup\{v\}$ is
OR-ed into the descendant bitset of every ancestor of $u$, and the dual update is
applied for ancestors; nodes outside the cone are untouched, so the work is
bounded by the size of the affected region rather than by $|V|$. Adding or
removing a variable allocates or frees a dense index and scrubs its incident
edges. A property flip on an edge's \textsf{active} flag is treated as an edge
insertion or deletion, and a flip of a variable's \textsf{latent} flag is a
property-only update that leaves the topology, and therefore the closures,
untouched.

\paragraph{When to fall back}
Not every edit patches cheaply. Edge \emph{deletion} can shrink ancestor and
descendant closures in non-local ways---a single removed edge may sever
reachability for many pairs---so when a deletion invalidates cached closures the
sound choice is to recompute the affected cone or, in the worst case, fall back to
a full recompilation. The system therefore treats incremental maintenance as an
optimization over the always-correct recompile: additive changes and property
flips patch in place, while topology-shrinking changes fall back when a bounded
patch would be unsound. \S\ref{sec:eval} measures the crossover---batch size and
edit type at which incremental maintenance beats recompilation
(Fig.~\ref{fig:incr}).

\NOTE{Implementation status (remove before CR; honest-gaps brief). The
transaction-listener plumbing---schema-change detection, the change-batch diff of
added/removed variables and edges, database-scoped registration/teardown---is
implemented. The post-commit hook currently applies the always-correct
\emph{fallback}: it invalidates the snapshot so the next query recompiles. The
bounded in-place patch described above (the OR-in over the affected cone) and the
pre-commit cycle/active/latent validators are specified but not yet wired; the
change batch is deliberately the context object through which they will be
applied. Frame R4 as ``designed, plumbing in place, patch pending'' until the
applyChange path and its eval land.}

%% Faithful to lifecycle/CausalGraphChangeListener.java, ChangeBatch.java,
%% CausalGraphLifecycle.java. See the status note on what is implemented vs.
%% designed.
%% ============================================================================
\section{Incremental Index Maintenance}\label{sec:incr}

A curated causal graph is edited continually---an edge is added as evidence
accrues, a variable is reclassified as latent, an edge is deactivated pending
review. Requirement R4 asks that the compiled index survive these edits without
a full recompilation after every change. This section describes how Causal4j
observes graph mutations and keeps the snapshot consistent.

\paragraph{Observing changes}
The plugin registers a transaction-event listener, scoped to the database for its
lifetime, that sees every committing write. Before a transaction commits, the
listener inspects the transaction data and, if it touches the causal schema,
distills the change into a compact \emph{change batch}: the ids of created and
deleted \textsf{Variable} nodes and the endpoint pairs of created and deleted
\textsf{CAUSES} edges. Transactions that touch no causal schema are rejected by a
cheap guard and cost nothing. The pre-commit hook is also the natural place to
\emph{validate}: because a well-formed causal model must remain acyclic, an edge
insertion that would close a cycle can be rejected here, aborting the offending
transaction rather than admitting an invalid model into the index. The batch is
handed to the post-commit hook, which applies it to the runtime once the write is
durable.

\paragraph{Incremental patch}
Given a change batch, the closures of \S\ref{sec:index:bitsets} can be maintained
locally rather than rebuilt. Inserting an edge $u\!\to\!v$ (having checked it
introduces no cycle) extends the CSR adjacency and then propagates closures over
only the \emph{affected cone}: $u$ and its ancestors gain $v$ and its descendants
as new descendants, and symmetrically $v$ and its descendants gain $u$ and its
ancestors as new ancestors. Concretely, $\textsf{descendantBits}[u]\cup\{v\}$ is
OR-ed into the descendant bitset of every ancestor of $u$, and the dual update is
applied for ancestors; nodes outside the cone are untouched, so the work is
bounded by the size of the affected region rather than by $|V|$. Adding or
removing a variable allocates or frees a dense index and scrubs its incident
edges. A property flip on an edge's \textsf{active} flag is treated as an edge
insertion or deletion, and a flip of a variable's \textsf{latent} flag is a
property-only update that leaves the topology, and therefore the closures,
untouched.

\paragraph{When to fall back}
Not every edit patches cheaply. Edge \emph{deletion} can shrink ancestor and
descendant closures in non-local ways---a single removed edge may sever
reachability for many pairs---so when a deletion invalidates cached closures the
sound choice is to recompute the affected cone or, in the worst case, fall back to
a full recompilation. The system therefore treats incremental maintenance as an
optimization over the always-correct recompile: additive changes and property
flips patch in place, while topology-shrinking changes fall back when a bounded
patch would be unsound. \S\ref{sec:eval} measures the crossover---batch size and
edit type at which incremental maintenance beats recompilation
(Fig.~\ref{fig:incr}).

\NOTE{Implementation status (remove before CR; honest-gaps brief). The
transaction-listener plumbing---schema-change detection, the change-batch diff of
added/removed variables and edges, database-scoped registration/teardown---is
implemented. The post-commit hook currently applies the always-correct
\emph{fallback}: it invalidates the snapshot so the next query recompiles. The
bounded in-place patch described above (the OR-in over the affected cone) and the
pre-commit cycle/active/latent validators are specified but not yet wired; the
change batch is deliberately the context object through which they will be
applied. Frame R4 as ``designed, plumbing in place, patch pending'' until the
applyChange path and its eval land.}


%% ----------------------------------------------------------------------------
%% ============================================================================
%% Section: Quantitative Layer and Expressiveness
%% Included from main.tex via %% ============================================================================
%% Section: Quantitative Layer and Expressiveness
%% Included from main.tex via %% ============================================================================
%% Section: Quantitative Layer and Expressiveness
%% Included from main.tex via \input{sections/quantitative}
%% Faithful to LinearGaussianInference.java, CounterfactualInference.java,
%% DefaultCausalGraphRuntime.java. KEEP SHORT -- expressiveness, not a contribution.
%% ============================================================================
\section{Quantitative Layer and Expressiveness}\label{sec:quant}

The same compiled snapshot that answers structural queries also carries optional
numeric parameters, which let Causal4j answer \emph{quantitative} causal
questions---interventions and counterfactuals---without leaving the DBMS. We
describe the layer briefly: it is a demonstration of the expressiveness the
indexed engine enables, not a contribution of the paper.

\paragraph{Linear-Gaussian model}
When a model is parameterized, each variable follows the structural equation
$V = \mu_{0,V} + \sum_{P\in\mathrm{pa}(V)} \beta_{P\to V}\,P + \varepsilon_V$ with
$\varepsilon_V\sim\mathcal{N}(0,\sigma^2_V)$, where the edge coefficients
$\beta$, intercepts $\mu_0$, and noise variances $\sigma^2$ are exactly the
\textsf{weight}, \textsf{intercept}, and \textsf{noise} properties read at compile
time (\S\ref{sec:compile}). The snapshot reports itself parameterized only when
some value departs from the structural-only defaults, so a purely qualitative DAG
incurs no numeric machinery.

\paragraph{Interventions and counterfactuals}
The do-operator $do(X{=}x)$ fixes $X$'s mean to $x$, sets its variance to zero,
and drops its incoming edges; the general recursion then handles it with no
special case. Interventional means propagate in topological order,
$\mu(V)=\mu_{0,V}+\sum_P \beta_{P\to V}\,\mu(P)$; interventional variances follow
from an incremental covariance recursion in the same order; and the total effect
of $T$ on $Y$ is obtained by unit-impulse propagation
($\partial T{=}1$, $\partial V=\sum_P \beta_{P\to V}\,\partial P$, read off at
$Y$), avoiding any path enumeration. Counterfactuals follow the standard
abduction--action--prediction recipe over the same equations: infer the exogenous
noise consistent with the observed factual world, apply the hypothesized
intervention, and re-propagate. Every one of these reuses the topological order
already on the snapshot, so the quantitative layer adds semantics, not a second
data structure.

\paragraph{Worked example}
On a parameterized version of the running example, \lstinline[style=causaldsl]|INTERVENE ON interventions = 1 COMPARE TO interventions = 0 RETURN EFFECT ON reported_cases|
returns the total effect along the front-door paths, while a
\lstinline[style=causaldsl]|COUNTERFACTUAL| statement answers what
\textsf{reported\_cases} would have been for a specific observed unit under a
different intervention. Because $do(\cdot)$ yields a derived model---the DAG with
$X$'s in-edges removed---these statements register a new causal world in the
model registry (\S\ref{sec:lang:registry}) that subsequent structural queries can
address like any other, closing the loop between the qualitative and quantitative
layers over a single stored graph.

%% Faithful to LinearGaussianInference.java, CounterfactualInference.java,
%% DefaultCausalGraphRuntime.java. KEEP SHORT -- expressiveness, not a contribution.
%% ============================================================================
\section{Quantitative Layer and Expressiveness}\label{sec:quant}

The same compiled snapshot that answers structural queries also carries optional
numeric parameters, which let Causal4j answer \emph{quantitative} causal
questions---interventions and counterfactuals---without leaving the DBMS. We
describe the layer briefly: it is a demonstration of the expressiveness the
indexed engine enables, not a contribution of the paper.

\paragraph{Linear-Gaussian model}
When a model is parameterized, each variable follows the structural equation
$V = \mu_{0,V} + \sum_{P\in\mathrm{pa}(V)} \beta_{P\to V}\,P + \varepsilon_V$ with
$\varepsilon_V\sim\mathcal{N}(0,\sigma^2_V)$, where the edge coefficients
$\beta$, intercepts $\mu_0$, and noise variances $\sigma^2$ are exactly the
\textsf{weight}, \textsf{intercept}, and \textsf{noise} properties read at compile
time (\S\ref{sec:compile}). The snapshot reports itself parameterized only when
some value departs from the structural-only defaults, so a purely qualitative DAG
incurs no numeric machinery.

\paragraph{Interventions and counterfactuals}
The do-operator $do(X{=}x)$ fixes $X$'s mean to $x$, sets its variance to zero,
and drops its incoming edges; the general recursion then handles it with no
special case. Interventional means propagate in topological order,
$\mu(V)=\mu_{0,V}+\sum_P \beta_{P\to V}\,\mu(P)$; interventional variances follow
from an incremental covariance recursion in the same order; and the total effect
of $T$ on $Y$ is obtained by unit-impulse propagation
($\partial T{=}1$, $\partial V=\sum_P \beta_{P\to V}\,\partial P$, read off at
$Y$), avoiding any path enumeration. Counterfactuals follow the standard
abduction--action--prediction recipe over the same equations: infer the exogenous
noise consistent with the observed factual world, apply the hypothesized
intervention, and re-propagate. Every one of these reuses the topological order
already on the snapshot, so the quantitative layer adds semantics, not a second
data structure.

\paragraph{Worked example}
On a parameterized version of the running example, \lstinline[style=causaldsl]|INTERVENE ON interventions = 1 COMPARE TO interventions = 0 RETURN EFFECT ON reported_cases|
returns the total effect along the front-door paths, while a
\lstinline[style=causaldsl]|COUNTERFACTUAL| statement answers what
\textsf{reported\_cases} would have been for a specific observed unit under a
different intervention. Because $do(\cdot)$ yields a derived model---the DAG with
$X$'s in-edges removed---these statements register a new causal world in the
model registry (\S\ref{sec:lang:registry}) that subsequent structural queries can
address like any other, closing the loop between the qualitative and quantitative
layers over a single stored graph.

%% Faithful to LinearGaussianInference.java, CounterfactualInference.java,
%% DefaultCausalGraphRuntime.java. KEEP SHORT -- expressiveness, not a contribution.
%% ============================================================================
\section{Quantitative Layer and Expressiveness}\label{sec:quant}

The same compiled snapshot that answers structural queries also carries optional
numeric parameters, which let Causal4j answer \emph{quantitative} causal
questions---interventions and counterfactuals---without leaving the DBMS. We
describe the layer briefly: it is a demonstration of the expressiveness the
indexed engine enables, not a contribution of the paper.

\paragraph{Linear-Gaussian model}
When a model is parameterized, each variable follows the structural equation
$V = \mu_{0,V} + \sum_{P\in\mathrm{pa}(V)} \beta_{P\to V}\,P + \varepsilon_V$ with
$\varepsilon_V\sim\mathcal{N}(0,\sigma^2_V)$, where the edge coefficients
$\beta$, intercepts $\mu_0$, and noise variances $\sigma^2$ are exactly the
\textsf{weight}, \textsf{intercept}, and \textsf{noise} properties read at compile
time (\S\ref{sec:compile}). The snapshot reports itself parameterized only when
some value departs from the structural-only defaults, so a purely qualitative DAG
incurs no numeric machinery.

\paragraph{Interventions and counterfactuals}
The do-operator $do(X{=}x)$ fixes $X$'s mean to $x$, sets its variance to zero,
and drops its incoming edges; the general recursion then handles it with no
special case. Interventional means propagate in topological order,
$\mu(V)=\mu_{0,V}+\sum_P \beta_{P\to V}\,\mu(P)$; interventional variances follow
from an incremental covariance recursion in the same order; and the total effect
of $T$ on $Y$ is obtained by unit-impulse propagation
($\partial T{=}1$, $\partial V=\sum_P \beta_{P\to V}\,\partial P$, read off at
$Y$), avoiding any path enumeration. Counterfactuals follow the standard
abduction--action--prediction recipe over the same equations: infer the exogenous
noise consistent with the observed factual world, apply the hypothesized
intervention, and re-propagate. Every one of these reuses the topological order
already on the snapshot, so the quantitative layer adds semantics, not a second
data structure.

\paragraph{Worked example}
On a parameterized version of the running example, \lstinline[style=causaldsl]|INTERVENE ON interventions = 1 COMPARE TO interventions = 0 RETURN EFFECT ON reported_cases|
returns the total effect along the front-door paths, while a
\lstinline[style=causaldsl]|COUNTERFACTUAL| statement answers what
\textsf{reported\_cases} would have been for a specific observed unit under a
different intervention. Because $do(\cdot)$ yields a derived model---the DAG with
$X$'s in-edges removed---these statements register a new causal world in the
model registry (\S\ref{sec:lang:registry}) that subsequent structural queries can
address like any other, closing the loop between the qualitative and quantitative
layers over a single stored graph.


%% ----------------------------------------------------------------------------
%% ============================================================================
%% Section: Evaluation
%% Included from main.tex via %% ============================================================================
%% Section: Evaluation
%% Included from main.tex via %% ============================================================================
%% Section: Evaluation
%% Included from main.tex via \input{sections/evaluation}
%% Owns Fig fig:latency and Fig fig:incr (referenced from \S7).
%% SKELETON: experimental design is written; all quantitative results are \TODO
%% pending measurement. See the status note.
%% ============================================================================
\section{Evaluation}\label{sec:eval}

\NOTE{Status: this section is the experimental \emph{design}. No results are
measured yet; every number is a \TODO. Do not submit without filling these. The
prose fixes the questions, baselines, datasets, and metrics so the harness has a
target; keep the ``pending'' framing until the artifact produces numbers.}

Our evaluation asks five questions. \textbf{RQ1 (query performance):} how do
Causal4j's structural and identification queries compare in latency and
throughput against alternatives? \textbf{RQ2 (index cost):} what does the
compiled index cost to build and to store, as a function of graph size and
density? \textbf{RQ3 (maintenance):} when does incremental maintenance beat full
recompilation? \textbf{RQ4 (scalability):} how does query latency grow with
$|V|$? \textbf{RQ5 (correctness):} do the answers agree with established causal
tools?

\paragraph{Baselines}
We compare four points along the ``how close to the DBMS'' axis: (a) an in-DBMS
map-of-sets adjacency with per-query DFS---the un-indexed baseline of
Table~\ref{tab:complexity}; (b) plain Cypher variable-length path queries
answering the same structural questions declaratively; (c) Neo4j GDS traversal
primitives~\cite{TODO_neo4j}; and (d) round-tripping the DAG to external causal
libraries---DoWhy~\cite{TODO_dowhy}, dagitty~\cite{TODO_dagitty},
pgmpy~\cite{TODO_pgmpy}, and causal-learn~\cite{TODO_causallearn}---which
includes export/import cost, the overhead our in-DBMS design removes.

\paragraph{Datasets}
We use standard Bayesian-network DAGs of increasing size (from small networks such
as \textsc{asia} and \textsc{sachs} through \textsc{alarm} and \textsc{hepar2} to
large ones such as \textsc{munin}), synthetic scale-free DAGs generated with a
preferential-attachment model to control $|V|$ and density independently, and the
COVID-19 running-example DAG~\cite{piccininni2020covid}. \TODO{Fix the exact
network list and the synthetic $|V|$ grid ($10^2$--$10^4$) and edge-density
settings; state the generator and seeds for reproducibility.}

\paragraph{Metrics and protocol}
For RQ1 we report median and tail latency and sustained throughput per operator
(ancestors/descendants, d-separation, adjustment-set enumeration). For RQ2, index
build time and resident memory vs.\ $|V|$ and density. For RQ3, wall-clock of an
incremental batch vs.\ a full recompile across batch sizes and edit types
(insert, delete, active/latent flip). For RQ4, latency vs.\ $|V|$ on a log--log
scale. For RQ5, exact agreement of d-separation verdicts and adjustment sets
against dagitty and causal-learn on every dataset. \TODO{State hardware, JVM
version and heap, Neo4j version, warmup and repeat counts, and report variance
(CI or IQR), per the reproducibility brief in the outline.}

\paragraph{Expected results (to be measured)}
On DAGs from \TODO{$V_{\min}$} to \TODO{$V_{\max}$} variables we expect Causal4j
to answer d-separation and adjustment-set queries \TODO{$N\times$} faster than
in-DBMS DFS and \TODO{$M\times$} faster than an external round-trip, to build its
index in \TODO{under X\,ms}, and to apply an incremental edge batch in
\TODO{Y} versus \TODO{Z} for a full recompile, while matching the reference tools
exactly on correctness (RQ5). Figures~\ref{fig:latency} and~\ref{fig:incr} report
RQ4 and RQ3 respectively.

\begin{figure}[t]\centering
  \fbox{\parbox{0.9\linewidth}{\centering \TODO{Fig.: query latency vs.\ $|V|$
  (log--log), one line per method (Causal4j index, DFS baseline, Cypher, GDS,
  external round-trip).}}}
  \caption{Query latency vs.\ graph size (RQ4). \emph{Pending measurement.}}
  \label{fig:latency}
\end{figure}

\begin{figure}[t]\centering
  \fbox{\parbox{0.9\linewidth}{\centering \TODO{Fig.: incremental update vs.\
  full recompile, by batch size and edit type; mark the crossover.}}}
  \caption{Incremental maintenance vs.\ recompilation (RQ3). \emph{Pending
  measurement.}}
  \label{fig:incr}
\end{figure}

%% Owns Fig fig:latency and Fig fig:incr (referenced from \S7).
%% SKELETON: experimental design is written; all quantitative results are \TODO
%% pending measurement. See the status note.
%% ============================================================================
\section{Evaluation}\label{sec:eval}

\NOTE{Status: this section is the experimental \emph{design}. No results are
measured yet; every number is a \TODO. Do not submit without filling these. The
prose fixes the questions, baselines, datasets, and metrics so the harness has a
target; keep the ``pending'' framing until the artifact produces numbers.}

Our evaluation asks five questions. \textbf{RQ1 (query performance):} how do
Causal4j's structural and identification queries compare in latency and
throughput against alternatives? \textbf{RQ2 (index cost):} what does the
compiled index cost to build and to store, as a function of graph size and
density? \textbf{RQ3 (maintenance):} when does incremental maintenance beat full
recompilation? \textbf{RQ4 (scalability):} how does query latency grow with
$|V|$? \textbf{RQ5 (correctness):} do the answers agree with established causal
tools?

\paragraph{Baselines}
We compare four points along the ``how close to the DBMS'' axis: (a) an in-DBMS
map-of-sets adjacency with per-query DFS---the un-indexed baseline of
Table~\ref{tab:complexity}; (b) plain Cypher variable-length path queries
answering the same structural questions declaratively; (c) Neo4j GDS traversal
primitives~\cite{TODO_neo4j}; and (d) round-tripping the DAG to external causal
libraries---DoWhy~\cite{TODO_dowhy}, dagitty~\cite{TODO_dagitty},
pgmpy~\cite{TODO_pgmpy}, and causal-learn~\cite{TODO_causallearn}---which
includes export/import cost, the overhead our in-DBMS design removes.

\paragraph{Datasets}
We use standard Bayesian-network DAGs of increasing size (from small networks such
as \textsc{asia} and \textsc{sachs} through \textsc{alarm} and \textsc{hepar2} to
large ones such as \textsc{munin}), synthetic scale-free DAGs generated with a
preferential-attachment model to control $|V|$ and density independently, and the
COVID-19 running-example DAG~\cite{piccininni2020covid}. \TODO{Fix the exact
network list and the synthetic $|V|$ grid ($10^2$--$10^4$) and edge-density
settings; state the generator and seeds for reproducibility.}

\paragraph{Metrics and protocol}
For RQ1 we report median and tail latency and sustained throughput per operator
(ancestors/descendants, d-separation, adjustment-set enumeration). For RQ2, index
build time and resident memory vs.\ $|V|$ and density. For RQ3, wall-clock of an
incremental batch vs.\ a full recompile across batch sizes and edit types
(insert, delete, active/latent flip). For RQ4, latency vs.\ $|V|$ on a log--log
scale. For RQ5, exact agreement of d-separation verdicts and adjustment sets
against dagitty and causal-learn on every dataset. \TODO{State hardware, JVM
version and heap, Neo4j version, warmup and repeat counts, and report variance
(CI or IQR), per the reproducibility brief in the outline.}

\paragraph{Expected results (to be measured)}
On DAGs from \TODO{$V_{\min}$} to \TODO{$V_{\max}$} variables we expect Causal4j
to answer d-separation and adjustment-set queries \TODO{$N\times$} faster than
in-DBMS DFS and \TODO{$M\times$} faster than an external round-trip, to build its
index in \TODO{under X\,ms}, and to apply an incremental edge batch in
\TODO{Y} versus \TODO{Z} for a full recompile, while matching the reference tools
exactly on correctness (RQ5). Figures~\ref{fig:latency} and~\ref{fig:incr} report
RQ4 and RQ3 respectively.

\begin{figure}[t]\centering
  \fbox{\parbox{0.9\linewidth}{\centering \TODO{Fig.: query latency vs.\ $|V|$
  (log--log), one line per method (Causal4j index, DFS baseline, Cypher, GDS,
  external round-trip).}}}
  \caption{Query latency vs.\ graph size (RQ4). \emph{Pending measurement.}}
  \label{fig:latency}
\end{figure}

\begin{figure}[t]\centering
  \fbox{\parbox{0.9\linewidth}{\centering \TODO{Fig.: incremental update vs.\
  full recompile, by batch size and edit type; mark the crossover.}}}
  \caption{Incremental maintenance vs.\ recompilation (RQ3). \emph{Pending
  measurement.}}
  \label{fig:incr}
\end{figure}

%% Owns Fig fig:latency and Fig fig:incr (referenced from \S7).
%% SKELETON: experimental design is written; all quantitative results are \TODO
%% pending measurement. See the status note.
%% ============================================================================
\section{Evaluation}\label{sec:eval}

\NOTE{Status: this section is the experimental \emph{design}. No results are
measured yet; every number is a \TODO. Do not submit without filling these. The
prose fixes the questions, baselines, datasets, and metrics so the harness has a
target; keep the ``pending'' framing until the artifact produces numbers.}

Our evaluation asks five questions. \textbf{RQ1 (query performance):} how do
Causal4j's structural and identification queries compare in latency and
throughput against alternatives? \textbf{RQ2 (index cost):} what does the
compiled index cost to build and to store, as a function of graph size and
density? \textbf{RQ3 (maintenance):} when does incremental maintenance beat full
recompilation? \textbf{RQ4 (scalability):} how does query latency grow with
$|V|$? \textbf{RQ5 (correctness):} do the answers agree with established causal
tools?

\paragraph{Baselines}
We compare four points along the ``how close to the DBMS'' axis: (a) an in-DBMS
map-of-sets adjacency with per-query DFS---the un-indexed baseline of
Table~\ref{tab:complexity}; (b) plain Cypher variable-length path queries
answering the same structural questions declaratively; (c) Neo4j GDS traversal
primitives~\cite{TODO_neo4j}; and (d) round-tripping the DAG to external causal
libraries---DoWhy~\cite{TODO_dowhy}, dagitty~\cite{TODO_dagitty},
pgmpy~\cite{TODO_pgmpy}, and causal-learn~\cite{TODO_causallearn}---which
includes export/import cost, the overhead our in-DBMS design removes.

\paragraph{Datasets}
We use standard Bayesian-network DAGs of increasing size (from small networks such
as \textsc{asia} and \textsc{sachs} through \textsc{alarm} and \textsc{hepar2} to
large ones such as \textsc{munin}), synthetic scale-free DAGs generated with a
preferential-attachment model to control $|V|$ and density independently, and the
COVID-19 running-example DAG~\cite{piccininni2020covid}. \TODO{Fix the exact
network list and the synthetic $|V|$ grid ($10^2$--$10^4$) and edge-density
settings; state the generator and seeds for reproducibility.}

\paragraph{Metrics and protocol}
For RQ1 we report median and tail latency and sustained throughput per operator
(ancestors/descendants, d-separation, adjustment-set enumeration). For RQ2, index
build time and resident memory vs.\ $|V|$ and density. For RQ3, wall-clock of an
incremental batch vs.\ a full recompile across batch sizes and edit types
(insert, delete, active/latent flip). For RQ4, latency vs.\ $|V|$ on a log--log
scale. For RQ5, exact agreement of d-separation verdicts and adjustment sets
against dagitty and causal-learn on every dataset. \TODO{State hardware, JVM
version and heap, Neo4j version, warmup and repeat counts, and report variance
(CI or IQR), per the reproducibility brief in the outline.}

\paragraph{Expected results (to be measured)}
On DAGs from \TODO{$V_{\min}$} to \TODO{$V_{\max}$} variables we expect Causal4j
to answer d-separation and adjustment-set queries \TODO{$N\times$} faster than
in-DBMS DFS and \TODO{$M\times$} faster than an external round-trip, to build its
index in \TODO{under X\,ms}, and to apply an incremental edge batch in
\TODO{Y} versus \TODO{Z} for a full recompile, while matching the reference tools
exactly on correctness (RQ5). Figures~\ref{fig:latency} and~\ref{fig:incr} report
RQ4 and RQ3 respectively.

\begin{figure}[t]\centering
  \fbox{\parbox{0.9\linewidth}{\centering \TODO{Fig.: query latency vs.\ $|V|$
  (log--log), one line per method (Causal4j index, DFS baseline, Cypher, GDS,
  external round-trip).}}}
  \caption{Query latency vs.\ graph size (RQ4). \emph{Pending measurement.}}
  \label{fig:latency}
\end{figure}

\begin{figure}[t]\centering
  \fbox{\parbox{0.9\linewidth}{\centering \TODO{Fig.: incremental update vs.\
  full recompile, by batch size and edit type; mark the crossover.}}}
  \caption{Incremental maintenance vs.\ recompilation (RQ3). \emph{Pending
  measurement.}}
  \label{fig:incr}
\end{figure}


%% ----------------------------------------------------------------------------
%% ============================================================================
%% Section: Related Work
%% Included from main.tex via %% ============================================================================
%% Section: Related Work
%% Included from main.tex via %% ============================================================================
%% Section: Related Work
%% Included from main.tex via \input{sections/related}
%% Citations use placeholder keys in refs.bib (TODO_*) -- verify before CR.
%% ============================================================================
\section{Related Work}\label{sec:related}

\paragraph{Reachability indexing}
Answering reachability without per-query traversal is a long line of database
work, from 2-hop labeling~\cite{TODO_2hop} through scalable indexes such as
GRAIL~\cite{TODO_grail}, pruned landmark labeling~\cite{TODO_pll}, and
topological-folding labels~\cite{TODO_tflabel}. Our ancestor/descendant bitsets
are a deliberately simple point in this space---a full transitive closure packed
into machine words---chosen because causal DAGs are comparatively small and the
queries above it (d-separation, adjustment enumeration) need whole ancestor and
descendant \emph{sets}, not just point-to-point reachability, which bitsets
deliver as one word-parallel operation. Positioning Causal4j in this lineage is
the point: we bring a classical reachability-indexing trade-off to a workload
(causal structural queries) that has not previously been treated as an indexed
DBMS workload.

\paragraph{Property-graph query languages}
Cypher and the emerging GQL/SQL-PGQ standards~\cite{TODO_gql} express graph
pattern matching and variable-length paths declaratively, and our language
compiles to Cypher plus procedures rather than replacing it. What Cypher does not
provide is conditioning-aware d-separation or adjustment-set enumeration as
first-class operators; these are not path patterns, and expressing them in Cypher
means client-side orchestration of many traversals. Causal4j adds the causal
operators while keeping the model queryable by ordinary Cypher.

\paragraph{Causal tooling}
Causal analysis is dominated by standalone libraries---DoWhy~\cite{TODO_dowhy},
dagitty~\cite{TODO_dagitty}, pgmpy~\cite{TODO_pgmpy},
causal-learn~\cite{TODO_causallearn}, and Ananke/Y0~\cite{TODO_ananke}---that
implement d-separation, the backdoor criterion, and identification over an
in-memory graph. None is DBMS-resident: each requires exporting the DAG, computes
in a separate process, and discards its representation, so none offers durable
shared access, in-database querying, or incremental maintenance under a changing
graph. We reuse their semantics (indeed we validate against them in
\S\ref{sec:eval}) but relocate the computation inside the store.

\paragraph{In-DBMS analytics and incremental maintenance}
Pushing analytics into the database via procedures and user-defined functions is
established practice; our contribution is the specific physical design and query
processing for causal operators, not the extension mechanism. Finally, keeping the
index consistent under updates connects to incremental view and index
maintenance~\cite{TODO_ivm}: our affected-cone patch of the reachability closures
(\S\ref{sec:incr}) is an instance of maintaining a materialized transitive
closure under edge insertion, with a principled fallback to recomputation for the
non-local case of deletion.

%% Citations use placeholder keys in refs.bib (TODO_*) -- verify before CR.
%% ============================================================================
\section{Related Work}\label{sec:related}

\paragraph{Reachability indexing}
Answering reachability without per-query traversal is a long line of database
work, from 2-hop labeling~\cite{TODO_2hop} through scalable indexes such as
GRAIL~\cite{TODO_grail}, pruned landmark labeling~\cite{TODO_pll}, and
topological-folding labels~\cite{TODO_tflabel}. Our ancestor/descendant bitsets
are a deliberately simple point in this space---a full transitive closure packed
into machine words---chosen because causal DAGs are comparatively small and the
queries above it (d-separation, adjustment enumeration) need whole ancestor and
descendant \emph{sets}, not just point-to-point reachability, which bitsets
deliver as one word-parallel operation. Positioning Causal4j in this lineage is
the point: we bring a classical reachability-indexing trade-off to a workload
(causal structural queries) that has not previously been treated as an indexed
DBMS workload.

\paragraph{Property-graph query languages}
Cypher and the emerging GQL/SQL-PGQ standards~\cite{TODO_gql} express graph
pattern matching and variable-length paths declaratively, and our language
compiles to Cypher plus procedures rather than replacing it. What Cypher does not
provide is conditioning-aware d-separation or adjustment-set enumeration as
first-class operators; these are not path patterns, and expressing them in Cypher
means client-side orchestration of many traversals. Causal4j adds the causal
operators while keeping the model queryable by ordinary Cypher.

\paragraph{Causal tooling}
Causal analysis is dominated by standalone libraries---DoWhy~\cite{TODO_dowhy},
dagitty~\cite{TODO_dagitty}, pgmpy~\cite{TODO_pgmpy},
causal-learn~\cite{TODO_causallearn}, and Ananke/Y0~\cite{TODO_ananke}---that
implement d-separation, the backdoor criterion, and identification over an
in-memory graph. None is DBMS-resident: each requires exporting the DAG, computes
in a separate process, and discards its representation, so none offers durable
shared access, in-database querying, or incremental maintenance under a changing
graph. We reuse their semantics (indeed we validate against them in
\S\ref{sec:eval}) but relocate the computation inside the store.

\paragraph{In-DBMS analytics and incremental maintenance}
Pushing analytics into the database via procedures and user-defined functions is
established practice; our contribution is the specific physical design and query
processing for causal operators, not the extension mechanism. Finally, keeping the
index consistent under updates connects to incremental view and index
maintenance~\cite{TODO_ivm}: our affected-cone patch of the reachability closures
(\S\ref{sec:incr}) is an instance of maintaining a materialized transitive
closure under edge insertion, with a principled fallback to recomputation for the
non-local case of deletion.

%% Citations use placeholder keys in refs.bib (TODO_*) -- verify before CR.
%% ============================================================================
\section{Related Work}\label{sec:related}

\paragraph{Reachability indexing}
Answering reachability without per-query traversal is a long line of database
work, from 2-hop labeling~\cite{TODO_2hop} through scalable indexes such as
GRAIL~\cite{TODO_grail}, pruned landmark labeling~\cite{TODO_pll}, and
topological-folding labels~\cite{TODO_tflabel}. Our ancestor/descendant bitsets
are a deliberately simple point in this space---a full transitive closure packed
into machine words---chosen because causal DAGs are comparatively small and the
queries above it (d-separation, adjustment enumeration) need whole ancestor and
descendant \emph{sets}, not just point-to-point reachability, which bitsets
deliver as one word-parallel operation. Positioning Causal4j in this lineage is
the point: we bring a classical reachability-indexing trade-off to a workload
(causal structural queries) that has not previously been treated as an indexed
DBMS workload.

\paragraph{Property-graph query languages}
Cypher and the emerging GQL/SQL-PGQ standards~\cite{TODO_gql} express graph
pattern matching and variable-length paths declaratively, and our language
compiles to Cypher plus procedures rather than replacing it. What Cypher does not
provide is conditioning-aware d-separation or adjustment-set enumeration as
first-class operators; these are not path patterns, and expressing them in Cypher
means client-side orchestration of many traversals. Causal4j adds the causal
operators while keeping the model queryable by ordinary Cypher.

\paragraph{Causal tooling}
Causal analysis is dominated by standalone libraries---DoWhy~\cite{TODO_dowhy},
dagitty~\cite{TODO_dagitty}, pgmpy~\cite{TODO_pgmpy},
causal-learn~\cite{TODO_causallearn}, and Ananke/Y0~\cite{TODO_ananke}---that
implement d-separation, the backdoor criterion, and identification over an
in-memory graph. None is DBMS-resident: each requires exporting the DAG, computes
in a separate process, and discards its representation, so none offers durable
shared access, in-database querying, or incremental maintenance under a changing
graph. We reuse their semantics (indeed we validate against them in
\S\ref{sec:eval}) but relocate the computation inside the store.

\paragraph{In-DBMS analytics and incremental maintenance}
Pushing analytics into the database via procedures and user-defined functions is
established practice; our contribution is the specific physical design and query
processing for causal operators, not the extension mechanism. Finally, keeping the
index consistent under updates connects to incremental view and index
maintenance~\cite{TODO_ivm}: our affected-cone patch of the reachability closures
(\S\ref{sec:incr}) is an instance of maintaining a materialized transitive
closure under edge insertion, with a principled fallback to recomputation for the
non-local case of deletion.


%% ----------------------------------------------------------------------------
%% ============================================================================
%% Section: Conclusion
%% Included from main.tex via %% ============================================================================
%% Section: Conclusion
%% Included from main.tex via %% ============================================================================
%% Section: Conclusion
%% Included from main.tex via \input{sections/conclusion}
%% ============================================================================
\section{Conclusion}\label{sec:concl}

Causal4j treats causal structural queries as a first-class, indexed workload
inside a property-graph DBMS. It contributes a declarative causal query language
that compiles to physical operators over the graph store, a compiled reachability
index---dense CSR adjacency with topological-layer pruning and precomputed
ancestor/descendant bitsets---that casts d-separation and adjustment-set
enumeration as bitset operations, and an incremental maintenance scheme that keeps
the index consistent as the graph is edited, all over a single authoritative
model that stays queryable by ordinary Cypher. Future work includes cost-based
selection between full-closure bitsets and layer-pruned traversal on large sparse
DAGs, an ADMG/latent-projection index for models with unmeasured confounding,
batched enumeration that shares work across the many d-separation tests an
identification query issues, and integration with the DBMS concurrency and
recovery model for multi-writer causal curation.

%% ============================================================================
\section{Conclusion}\label{sec:concl}

Causal4j treats causal structural queries as a first-class, indexed workload
inside a property-graph DBMS. It contributes a declarative causal query language
that compiles to physical operators over the graph store, a compiled reachability
index---dense CSR adjacency with topological-layer pruning and precomputed
ancestor/descendant bitsets---that casts d-separation and adjustment-set
enumeration as bitset operations, and an incremental maintenance scheme that keeps
the index consistent as the graph is edited, all over a single authoritative
model that stays queryable by ordinary Cypher. Future work includes cost-based
selection between full-closure bitsets and layer-pruned traversal on large sparse
DAGs, an ADMG/latent-projection index for models with unmeasured confounding,
batched enumeration that shares work across the many d-separation tests an
identification query issues, and integration with the DBMS concurrency and
recovery model for multi-writer causal curation.

%% ============================================================================
\section{Conclusion}\label{sec:concl}

Causal4j treats causal structural queries as a first-class, indexed workload
inside a property-graph DBMS. It contributes a declarative causal query language
that compiles to physical operators over the graph store, a compiled reachability
index---dense CSR adjacency with topological-layer pruning and precomputed
ancestor/descendant bitsets---that casts d-separation and adjustment-set
enumeration as bitset operations, and an incremental maintenance scheme that keeps
the index consistent as the graph is edited, all over a single authoritative
model that stays queryable by ordinary Cypher. Future work includes cost-based
selection between full-closure bitsets and layer-pruned traversal on large sparse
DAGs, an ADMG/latent-projection index for models with unmeasured confounding,
batched enumeration that shares work across the many d-separation tests an
identification query issues, and integration with the DBMS concurrency and
recovery model for multi-writer causal curation.


%% ----------------------------------------------------------------------------
%% The AI-content acknowledgement does NOT count toward the 12-page limit.
\section*{Acknowledgment on AI-Generated Content}
\NOTE{Verify against the ICDE 2027 call before CR; adjust to match actual use.}
The authors used an AI assistant (Anthropic Claude) to help draft and edit the
prose of this paper and to generate the \TeX{} of its figures, working from the
authors' outline, code base, and experimental design. All technical content---the
system design, algorithms, and claims---originates with the authors, who
reviewed, corrected, and take responsibility for the final text. No AI system was
used to produce experimental results or measurements.

\bibliographystyle{IEEEtran}
\bibliography{refs}

\end{document}
